
%
%
%
\section{Introduction}\label{sec:intro}

For $n\in\mathbb{N}$, the problem of generating permutations of 
$[n]=\{1,2\dots,n\}$ at random is foundational in the history of 
computer science~\cite{knuth2}. Markov chains for sampling permutations 
arise in a variety of contexts, including self-organizing 
lists~\cite{selforganizingsurvey, rivest}, card shuffling~\cite{bd, 
wilson}, and search engines~\cite{search}. The spectral gap of a Markov 
chain provides a measure of the rate of convergence to stationarity, 
which is crucial to the efficiency of Markov chain algorithms for 
sampling.

Suppose we are given %
%
an array  %
of input probabilities $\cP = \{p_{i,j}\}$ 
%
%
%
$(1 \leq i, j \leq n)$   %
with $p_{i, j} = 1-p_{j, i}$.  A natural 
Markov chain $\mn$ over permutations operates by uniformly choosing a 
pair of adjacent elements, $i$ and~$j$, and putting~$i$ ahead of~$j$ 
with probability $p_{i,j}$ and~$j$ ahead of~$i$ with probability 
$p_{j,i}$, %
%
regardless   %
of their current ordering.  We call $\mn$ the 
\emph{nearest-neighbor} transposition chain.

The Markov chain $\mn$ was among the first Markov chains studied in 
terms of its computational efficiency for sampling \cite{ald, diasha, 
dsc}. Its spectral gap has been studied extensively, both in 
the uniform and
in the  %
biased settings \cite{BBHM05,bmrs,DSC93b, diasha, wilson}. 

%
%
A central question is, %
under what conditions %
is there a lower bound on the spectral gap that is
is an inverse polynomial in $n$.
%
Such a lower bound    %
implies a polynomial
%
bound 
on the \emph{mixing time}, or the time until the chain will be 
``close'' to its stationary distribution. We say $\cP$ is 
\emph{positively biased} if $p_{i,j}\geq 1/2$ for all $i<j$.
Without the %
%
positive bias  %
restriction, it is easy to find examples where the mixing time of $\mn$ is
not polynomial (see, \eg,~\cite{bmrs}); 
however, even with this restriction the mixing time can still not be polynomial, as demonstrated in~\cite{bmrs}.  In~\cite{bmrs}, the authors give an example of positively biased probabilities $\cP$ for which they prove the mixing time of $\mn$ is exponential in $n.$

In 2003, %
Fill~\cite{F03}  %
introduced the following monotonicity 
conditions: $p_{i,j+1}\geq p_{i,j}$ for all $1\leq i<j\leq n-1$ and 
$p_{i-1,j}\geq p_{i,j}$ for all $2\leq i<j\leq n$.  He conjectured that  
if $\cP$ is positively biased and monotone, then the spectral gap of 
$\mn$ is  %
at least  %
an inverse polynomial in $n$, and moreover that the smallest 
spectral gap is attained when each $p_{i,j}=1/2$ and the stationary 
distribution is uniform. 

Despite significant work on this subject, Fill's conjecture remains 
mostly open after more than 15 years. In the uniform case, Wilson~\cite{wilson} produced a 
clever path coupling argument showing the mixing 
time is $\Theta(n^3\log n)$, with upper and lower bounds within a factor of 2. 
Fill's conjecture states that this should be the worst case mixing time among all 
arrays $\cP$ which are positively biased and monotone. 
%
%
%
%
%
Various papers~\cite{BBHM05, bmrs} have identified 
different classes of
arrays   %
$\cP$ for which $\mn$ may cleverly be viewed as 
the direct product of simpler independent Markov chains, and thus may 
be analyzed easily in terms of those chains. In~\cite{BBHM05}, the 
authors proved a bound of $O(n^2)$ on the mixing time of $\mn$ in the 
case that $p_{i,j} = p$ for all $i<j$, and in~\cite{bmrs} the authors 
considered the case that $p_{i,j}$ depends only on the smaller of $i$ 
and $j$.

Among the positively biased, monotone %
%
arrays $\cP$  %
that have proven to 
be challenging to analyze are those arising in the context of 
self-organizing lists, where each element $i$ has a frequency $w_i$ of 
being requested, and then moved ahead
by  %
one in the list; in this case, 
$p_{i,j}=w_i/(w_i+w_j)$.  The Markov chain $\mn$ was termed a 
``gladiator chain'' in this case by Haddadan and Winkler~\cite{HW16}. 

\begin{figure}[htb]
  \centering
       \begin{minipage}[c]{.5\textwidth}
     \centering
     {\footnotesize
     A partition of $[7]$:
     \begin{tabular}{rl}
     $\mathcal{C}_1$ & $= \{{\color{red} 1,2,3}\}$\\
      $\mathcal{C}_2$ & $= \{{\color{blue} 4,5}\}$\\
      $\mathcal{C}_3$ & $= \{{\color{green} 6}\}$\\
       $\mathcal{C}_4$ & $= \{{\color{black} 7}\}$\\
     
     \end{tabular}}
     \end{minipage}\hfill
          \begin{minipage}[c]{.5\textwidth}
     \centering
           \newcommand{\one}{{\color{red}1}}
           \newcommand{\two}{{\color{blue}2}}
           \newcommand{\three}{{\color{green}3}}

          \newcommand{\four}{{\color{black}4}}

    {\color{red} 3}{\color{black}7}{\color{green}6}{\color{blue}5}{\color{red}2}{\color{blue}4}{\color{red}1} $\rightarrow \one \four \three \two \one\two\one$\\
    
       \end{minipage}\hfill
      \caption{An example of a 4-class permutation and corresponding 4-particle process, with $n=7$.}
\label{fig:Particle}
\end{figure}

It is instructive to consider \emph{$k$-classes}~\cite{MS}, where $[n]$ 
is partitioned into $k$ classes and particles from class $i$ and class 
$j$ interact with a fixed probability
%
%
$\pbar_{i,j}\,$.\footnote{The bar in   %
the notation indicates that we have re-indexed the probabilities by 
class.} %
When $k=n$, this captures the general setting. In this 
context, $\mn$ can be seen as having dual duties: \emph{whisking}, 
which uniformly mixes particles of the same type\footnote{We use the 
terms ``type'' and ``class'' interchangeably.}, and \emph{sifting}, 
which mixes particles of different types %
-- that is, which changes relative orders of particles of different types~\cite{HW16}. %
%
As the mixing properties of the uniform chain are well-understood, it is 
sufficient to analyze the sifting operation in
%
%
isolation \cite{HW16, MS}.  %
By discarding moves between particles in the same class, we are 
left with a linear $k$-particle process that maintains elements within 
each particle class in fixed relative orders (and therefore we may drop 
their individual labels and re-index, identifying all elements from 
class $i$ with the label $i$, as is done in \expref{Figure}{fig:Particle}).

In this paper, we use the version of the $k$-particle process 
introduced in~\cite{MS}, called $\mk$, which is also allowed to make 
certain non-adjacent transpositions---it may swap $i$ and $j$ if all 
elements between them are smaller than both $i$ and $j$.  This 
simplifies our analysis, and as in~\cite{MS}, we compensate by using 
comparison techniques~\cite{dsc, RT98} when evaluating the spectral gap 
of $\mn$.
 
The \emph{bias} towards having a particle of type $i$ ahead of a 
particle of type $j$ is $\qbar_{i,j} = \pbar_{i,j}/\pbar_{j,i}$.  Let 
$q=\min_{i,j}\qbar_{i,j}$ be the minimum bias. We say that the bias is 
\emph{bounded} if $q$ is bounded below by a constant greater than 1. 
After a series of papers~\cite{HW16,MS}, it was shown that the mixing time
of $\mk$ was $O(n^{2k+4})$ (and the spectral 
gap is at least $\Omega(n^{-2(k-1)})$) 
which $\cP$ is positively biased, monotone, and bounded. %
%
%
%
%
%
%
These results apply to the 
gladiator chain (self-organizing lists) with $k$ distinct frequencies.  
In fact, the result in~\cite{MS} requires only \emph{weak 
monotonicity}, and not full monotonicity.  Weak monotonicity for 
$k$-classes is defined as follows.

\begin{definition}[\cite{bmrs}]\label{mono} If $\cP$ forms a $k$-class, 
then $\cP$ is weakly monotonic if properties 1 and either 2 or 3 are satisfied.
\begin{enumerate}
\item $\pbar_{i,j} \geq 1/2$ for all $1 \leq i < j \leq k,$ and
\item $\pbar_{i,j+1} \geq \pbar_{i,j}$ for all $1 \leq i < j \leq k-1$ or
\item $\pbar_{i-1,j} \geq \pbar_{i,j}$ for all $2 \leq i < j \leq k.$
\end{enumerate}
\end{definition}

The aforementioned results are based on a natural decomposition of 
$\mk$ into simpler chains, but not as a direct product. To get a bound 
on the overall spectral gap, the authors of~\cite{MS} used the 
decomposition theorem of~\cite{MR00}, which bounds the spectral gap of 
a Markov chain in terms of the spectral gaps of the simpler Markov 
chains. Unfortunately, one incurs an extra factor of $n^{-2}$ each time 
it is applied in this setting, and in~\cite{MS} it is applied 
iteratively $k-2$ times. Thus, this produced a bound
%
of $\Omega(n^{-2(k-1)})$ on the spectral gap,
which is an inverse polynomial only for constant $k$.

To make this iterated decomposition scheme work for larger $k$ requires 
a stronger decomposition theorem, and that is the main focus of the 
present paper. We introduce a new decomposition theorem that allows us 
to achieve nearly optimal bounds of $\Omega(n^{-2})$ on the spectral 
gap of $\mk$ for bounded \kPs, as long as the number of particles of 
each type is at least $C_q\log n$ (where $C_q$ is a constant depending 
on the minimum bias $q$; \ie,  not depending on $n$).  We believe this 
new decomposition theorem is of independent interest and will have 
other applications.  

\subsection{The decomposition method}

The decomposition method was first introduced by Madras and 
Randall~\cite{madr}, and has been subsequently used and modified to 
produce the first polynomial time bounds on the spectral gaps of many 
interesting Markov chains \cite{CDF00,Dest03,Ding2010, EMT18, 
GHS08,HW16, jstv04,mh15,MR00,mrs,PS17, R01}. Suppose $\m$ is a finite, 
ergodic Markov chain that is reversible and has stationary 
distribution $\pi$. Let
%
$\Omega=\bigcup_{i=1}^{\rhat}\Omegahat_i$  %
be a partition of the state space of $\m$ and $\gammahat_i$ 
be the spectral gap of $\m$ restricted to $\Omegahat_i$. The disjoint 
decomposition theorem of~\cite{MR00} states that the spectral gap 
$\gamma$ of $\m$ satisfies $\gamma \geq \frac{1}{2}\gmin\gammabar$, 
where $\gmin=\min_i \gammahat_i$ and $\gammabar$ is the spectral gap of 
a certain \emph{projection} chain over states $[r]$. 

There has been significant effort towards improving the decomposition 
technique by providing stronger bounds in special 
cases \cite{Dest03,EMT18,GHS08,jstv04,mh15,MR00,PS17,R01}. While 
$\gamma$ may indeed be on the order of $\gmin\gammabar$---one example 
is the random walk on the path, decomposed into two
%
shorter  %
paths---there are instances in which it may instead satisfy the much 
larger bound $\gamma \geq c \min\{\gmin, \gammabar\}$, for some 
constant $c$. The simplest such example is the direct product of two 
independent Markov chains~\cite{bmrs,EMT18}; in this case, $c=1$. 

Tight bounds are especially important when applying the decomposition 
method iteratively, as was done in~\cite{MS}. At each level of the 
induction, $\gammabar=\Theta(n^{-2})$, so the original bound 
of~\cite{MR00} yields $\gamma=\Omega(n^{-2(n-1)})$ for the final 
iteration.   Even a bound of the form $\gamma \geq c \min\{\gmin, 
\gammabar\}$, such as the one in~\cite{PS17}, would introduce a factor 
of $c$ for each application, and would thus yield a bound that is an 
inverse exponential in $n$ if $c<1$ is constant. The bounds 
in~\cite{jstv04} are iterable in some cases, but $\mk$ does not satisfy 
those conditions.

\subsection{Our results}

In this paper, we develop a new decomposition framework that yields 
iterable bounds for a new class of Markov chains. Among our results, we 
present a complementary decomposition theorem, which achieves a tight 
bound on $\gamma$ without appealing to a bound on the gap $\gammabar$ 
of the projection chain, but rather the minimum gap $\gmintilde$ of 
certain \emph{complementary restrictions} 
$\Ptilde_1,\Ptilde_2,\ldots,\Ptilde_{\rtilde}$. We first consider the 
simple setting where the state space $\Omega$ can be seen as a product 
space, \ie,  $\Omega=\Omega_1\times\Omega_2$. In other words, for each 
$a\in\Omega_1$ and each $b\in\Omega_2$, there is a unique 
$\sigma=(a,b)\in \Omega$. This setting is similar to the direct product 
of independent Markov chains, but the transition probabilities are not 
necessarily independent. We define a restriction chain $P_a$ for each 
$a\in\Omega_1$ that fixes $a$ and operates only on the second 
coordinate.  Similarly, we define a complementary restriction chain 
$\Ptilde_b,$ which fixes $b$ and operates only on the first coordinate. 
Recall~$\pi$ is the stationary distribution of $\m$. We write $\pi(a) = 
\sum_{b\in \Omega_2}\pi(a,b)$ and $\pi(b) = \sum_{a\in 
\Omega_1}\pi(a,b)$.  Define \[r(a,b) = \frac{\pi(a,b)}{\pi(a)\pi(b)}\,.\] 
The function $r(a,b)$ allows us to capture the degree of dependence 
between $a$ and $b$.  Let
\begin{equation}\label{eq:eps}
\epsilon = \sum_{(a,b)\in\Omega}\pi(a,b)\left(\sqrt{r(a,b)} -
1/\sqrt{r(a,b)}\right)^2.
\end{equation}
We say a decomposition satisfying the equality above is \emph{$\epsilon$-orthogonal}. 
%
%
%
%
%

\begin{theorem}\label{thm:productspaces}
For any $\epsilon$-orthogonal decomposition of Markov chain $\m$ on 
product space $\Omega$,
\[ \gamma(\m) \geq \min\{\gminhat, \gmintilde\}\left(1 -\sqrt{\epsilon}\right)^2\,.\]
\end{theorem}

\noindent This bound can be iterated $t$ times with 
only a constant overhead, as long as $\sqrt{\epsilon}\leq 1/t$.
We note that parts of the proof of this theorem are similar
to a ``multi-decomposition'' result of Destainville~\cite{Dest03},
which we discuss in \expref{Section}{sec:generalizations}.  There we also
present several generalizations of \expref{Theorem}{thm:productspaces}, 
which apply beyond the product space setting.

Favorably, analysis of $\epsilon$-orthogonality requires only a 
comparison between two stationary distributions and not an analysis of 
the dynamics of any Markov chain.  For example, when $\m$ is a direct 
product of independent Markov chains, we have that $r(a,b)=1$ for all 
pairs $a\in\Omega_1$ and $b\in\Omega_2$ and the decomposition is 
$0$-orthogonal, leading to the bound $\gamma\geq 
\min\{\gminhat,\gmintilde\}$, as expected.  Notice, however, that we do 
not actually require a such a strong pointwise bound on $r(a,b)$.  The 
notion of $\epsilon$-orthogonality captures the \emph{average} value of 
$r(a,b)$, and allows us to achieve tight bounds on $\gamma$ even when 
the constituent Markov chains are only \emph{nearly} independent.  
Indeed, it is possible to prove %
$\epsilon$-orthogonality for very small $\epsilon$ %
even if 
$r(a,b)$ is far from 1 for pathological pairs $a$ and $b$, as long as 
it is close to 1 on average. Importantly, this holds even though the 
elements in this ``bad'' space are visited polynomially often.  

Armed with our new decomposition theorem, we use the iterated 
decomposition in~\cite{MS} to achieve nearly optimal bounds on the 
spectral gap of the $k$-particle process~$\mk$. We prove that at each
level of the iteration, %
the value of $\epsilon$ in the decomposition satisfies $\epsilon\leq 1/n^2$. %
%
Thus, we may apply \expref{Theorem}{thm:productspaces} iteratively $k\leq 
n$ times to get a bound of $\Omega(n^{-2})$ on the spectral gap 
$\gamma(\mk)$.  This bound is optimal up to constants. More formally, 
let $\psize = \max\left\{\frac{\log 
(6n^2)+\log((1+q)/(q-1))}{\log((1+q)/2)},  
\frac{\log^2(14)}{\log^2(2q/(1+q))}\right\}$, %
and let 
$c_i$ denote the number of particles of type $i$, for $1\leq i\leq k$.  
Note that $\psize := C_q \log n$  where $C_q$ depends only on $q.$
We prove the following.

\newtheorem*{thm:mk}{\expref{Theorem}{pprocess}}

\begin{theorem}\label{pprocess}
%
 %
%
If the probability array $\cP$ is   %
weakly monotonic and %
forms  %
a bounded $k$-class
with  $c_i\geq 2\psize$ for all $1\leq i \leq k$, 
then the spectral gap $\gamma$ of the chain $\mk$ satisfies
\abovedisplayskip=3pt
\belowdisplayskip=3pt
$\gamma =\Omega(n^{-2}).$
\end{theorem}

\begin{figure}[]
  \centering
    \newcommand{\one}{{\color{gray}1}}
     {\color{red} 4\one 3\one \one \underline{2}4\one 3\one 42}\\
    \caption{At level 2 of the decomposition, all particles of type 1 are in fixed positions, and the underlined particle 2 may swap with the 3 to its left or the 4 immediately to its right.}
\label{fig:SwapAcross}
\end{figure}

The iterated decomposition works as follows. Let $\mathcal{C}_{i}$ be 
the set of particles of type $i$, $\mathcal{C}_{<i}$ the set of 
particles with type less than $i$, and $\mathcal{C}_{>i}$ the set of 
particles with type bigger than $i$. At the $i$th iteration, all 
particles in $\mathcal{C}_{<i}$ are in fixed positions, and those in 
$\mathcal{C}_{i}\cup \mathcal{C}_{>i}$ are allowed to swap across them
-- that is, we allow non-adjacent transpositions of, for example, a particle 
of type $j\geq i$ and a particle of type $j'\geq i$, as long as all particles 
between them in the permutation are of type less than $i$. %
%
%
%
See \expref{Figure}{fig:SwapAcross}
for an example.  %
This decomposition is designed to 
exploit the hypothesis that the movement of the particles in 
$\mathcal{C}_{i}$ is nearly independent of the relative order of the 
particles in $\mathcal{C}_{>i}$.  The tool of $\epsilon$-orthogonality 
allows us to make this precise.  We define the complementary 
restriction chains to contain the moves involving only particles in 
$\mathcal{C}_{>i}$, and we define the restriction chains to contain 
moves between particles in $\mathcal{C}_{i}$ and particles in 
$\mathcal{C}_{>i}$. We prove that this decomposition is 
%
$\epsilon-$orthogonal where $\epsilon\leq 1/n^2$ %
when $|\mathcal{C}_{i+1}| = c_{i+1}$ is large 
enough.  Indeed, the highest probability configuration is the one in 
which particles are sorted by class, with all particles of smaller type 
appearing before particles of larger type.  Thus, having 
$|\mathcal{C}_{i+1}|$ large ensures that a typical configuration will 
not have a particle in $\mathcal{C}_{i}$ after a particle in 
$\mathcal{C}_{>(i+1)}$, as this requires many transpositions from the 
highest probability configuration, and each costs a factor of at least 
the minimum bias $q$; see \expref{Figure}{fig:typical}.  Note that such 
``bad'' configurations are polynomially suppressed 
(meaning, their total probability is bounded by an inverse polynomial in $n$), %
but not exponentially suppressed
(bounded by an inverse exponential in $n$).  %
%
%

\begin{figure}[]
  \centering
           \newcommand{\one}{{\color{gray}1}}
         {\color{red} \one 223\one 23\one \one 2\one 334\one 3\one 43444}\\
      \caption{A ``typical'' configuration at level 2 has no 2's appearing after any $j\geq 4$.}
\label{fig:typical}
\end{figure}

We use \expref{Theorem}{pprocess} and comparison techniques to prove the following result for $\mn$.

\newtheorem*{mnGap}{\expref{Theorem}{thm:mnGap}}
\begin{theorem}\label{thm:mnGap} If the
%
probability array $\cP$ is  %
weakly monotonic and
%
forms  %
a bounded \kP\, with at least~$2\psize$ particles 
 in each class, then $\Gap(\mn)=\Omega(n^{-7})$.
\end{theorem}
 
\noindent This is the first inverse-polynomial %
lower   %
bound on the spectral 
gap of $\mn$ for monotone bounded $k$-classes for $k$ as large as 
$\Theta(n/\log n)$.  This is a significant improvement over the 
previous result which was inverse polynomial only for constant 
$k$~\cite{MR00}.  \expref{Theorem}{thm:mnGap} also leads to a bound of 
$O(n^9)$ on the mixing time of $\mn$ under the same conditions.

\subsection{Outline}

The layout of the paper is as follows.  In \expref{Section}{sec:prelim}, we 
begin with some notation and terminology.  In 
\expref{Section}{sec:decompThm}, we give details on 
\expref{Theorem}{thm:productspaces} and illustrate its use by applying it 
to the one-dimensional Ising model. In 
\expref{Section}{sec:generalizations}, we generalize the notion of 
$\epsilon$-orthogonality to non-product spaces and present a more 
general complementary decomposition theorem that applies to all 
$\epsilon$-orthogonal decompositions. We also present a classical 
decomposition theorem that generalizes some results of~\cite{jstv04} 
and show how our decomposition theorems relate to previous results. 
Finally, we give a brief summary of the proof techniques for the 
decomposition theorems.  \expref{Section}{sec:decompProofs} 
and~\expref{Section}{sec:classical} present the proofs of the complementary and 
classical decomposition theorems, respectively. In 
\expref{Section}{sec:perm}, we apply \expref{Theorem}{thm:productspaces} to the 
biased permutation problem. 


\section{Preliminaries}\label{sec:prelim}



We assume $\m$ is an ergodic Markov chain over a finite state space 
$\Omega$ with transition matrix $P$. We also assume $\m$ is reversible 
and has stationary distribution $\pi$; that is, it 
satisfies the \emph{detailed balance} condition: for all 
$\sigma,\tau\in\Omega$, 
$\pi(\sigma)P(\sigma,\tau)=\pi(\tau)P(\tau,\sigma)$.  Notationally, 
we write $\pi(S) = \sum_{\elemone\in S} \pi(\elemone)$ for any set 
$S\subseteq \Omega$.

Let $\Omega = \cup_{i=1}^{r} \Omega_i$ be a partition of the state 
space into $r$ pieces.  For each $i \in [r]$, define $P_{i} = 
\transmat(\Omega_i)$ as the restriction of $\transmat$ to $\Omega_i$ 
which rejects moves that leave $\Omega_i$.  
Formally, $P_{i}$ is defined as 
follows: if $\elemone \not= \elemtwo$ and $\elemone,\elemtwo \in 
\Omega_i$ then $P_i(\elemone,\elemtwo)=\transmat(\elemone,\elemtwo)$; 
if $\elemone \in \Omega_i$ then 
$P_i(\elemone,\elemone)=1-\sum_{\elemtwo \in \Omega_i, \elemtwo \not= 
  \elemone} P_i(\elemone,\elemtwo)$.  The Markov chain $\m_i$ with transition matrix $P_i$ is called a \emph{restriction} Markov chain, and its state space is $\Omega_i$.
Let $\pi_i$ be the normalized restriction of 
$\pi$ to $\Omega_i$; \ie,  $\pi_i(S)=\pi(S\cap\Omega_i)/\pi(\Omega_i)$ 
for any $S \subseteq \Omega$.  The chain $\m_i$ is ergodic, 
reversible, and has stationary distribution $\pi_i$.

Classical decomposition theorems use the so-called \emph{projection} chain $\bar{\m}$
with transition matrix $\Pbar$ on the state space $[\rhat]$ defined by 
$\Pbar(i,j)=\pi(\Omegahat_i)^{-1}\sum_{{\elemone\in\Omegahat_i,}{\elemtwo
\in \Omegahat_j}} \pi(\elemone)\transmat(\elemone,\elemtwo)$.
The stationary distribution $\pibar$ of $\bar{\m}$ satisfies $\pibar(i) := 
\pi(\Omegahat_i)$. 

In contrast, for our complementary decomposition theorems,
we will decompose $\transmat$ in a slightly different way.
Define $\Phat$ as the block diagonal $|\Omega|\times |\Omega|$ matrix 
with the $P_i$ matrices along the diagonal; \ie,  $\Phat$ is 
obtained from $\transmat$ by rejecting moves between different parts of 
the partition.
Define $\Ptilde$ to be the transition matrix of the 
Markov chain defined by rejecting moves from $\elemone$ to $\elemtwo$ 
if $\elemone$ and $\elemtwo$ are within the same $\Omegahat_i$.  
The matrix~$\Ptilde$ 
defines a \emph{complementary} partition $\Omega = \cup_{j=1}^{\rtilde} 
\Omegatilde_j$, where each $\Omegatilde_j$ is a maximal subset of 
$\Omega$ that is connected 
by $\Ptilde$. For each $j\in [\rtilde]$, define the \emph{complementary 
restriction} $\Ptilde_{j} = \transmat(\Omegatilde_j)$ as the 
restriction of $\transmat$ to $\Omegatilde_j$ which rejects moves that 
leave $\Omegatilde_j$. The complementary restriction $\Ptilde_j$ is 
also ergodic, reversible, and its stationary distribution is the normalized restriction 
of $\pi$ to $\Omegatilde_j$, which we call $\pitilde_j$.
Notice that 
the complementary restrictions are defined by the decomposition 
$P_1,P_2,\ldots, P_r$.

Observe $(\Phat+\Ptilde)(\elemone,\elemtwo)=\transmat(\elemone,\elemtwo)$ 
unless $\elemone=\elemtwo$, and 
$(\Phat+\Ptilde)(\elemone,\elemone)=\transmat(\elemone,\elemone)+1,$ 
since each move of $\transmat$ gets rejected in exactly one of $\Phat$ 
or $\Ptilde$ (and of course the probability of transitioning from a 
state is 1).  Therefore, we have $\transmat = 
\Phat+\Ptilde-I_{|\Omega|}$, where we write $I_n$ to 
mean the $n \times n$ identity matrix. 

The efficiency of a Markov chain $\m$ is a function of its spectral 
gap, denoted $\gamma(\m)$, which is defined as the difference of 1 and 
the second largest eigenvalue of its transition matrix. Letting 
$\gammahat_i = \gamma(\Omega_i)$ and $\gammatilde_j = 
\gamma(\Omegatilde_j)$, the complementary decomposition theorem, 
\expref{Theorem}{thm:productspaces}, is proven by analyzing the spectral 
gaps $\gminhat=\min_i \gammahat_i$ and $\gmintilde=\min_j 
\gammatilde_j$.  Note that if some restriction or complementary 
restriction has a single element, its spectral gap is taken to be $1$.

\section{Introductory examples}\label{sec:decompThm} 

In this section, we show how to apply our new complementary
decomposition theorem by considering a few simple examples.  Recall
$r(a,b) = \pi(a,b)/(\pi(a)\pi(b))$ and
\vspace{-.1in}
\[\epsilon = \sum_{(a,b)\in\Omega}\pi(a,b)\left(\sqrt{r(a,b)} -
1/\sqrt{r(a,b)}\right)^2\,.\]

\vspace{-.15in}

\noindent \expref{Theorem}{thm:productspaces} states that the spectral gap 
$\gamma$ of $\m$ satisfies $\gamma \geq \min\{\gminhat, 
\gmintilde\}\left(1 -\sqrt{\epsilon}\right)^2$. 
A simple application of \expref{Theorem}{thm:productspaces} is to a Markov 
chain $\m$ that is the direct product of two Markov chains $\m_1$ and $ 
\m_2$.  It is easy to see that $r(a,b)=1$ for all $a,b$, and so this proves
$\gamma(\m) = \min\{\gamma(\m_1), \gamma(\m_2)\}$.  By 
iterating on $\gmintilde$, we can immediately prove the following well-known result. 

\begin{corollary}\label{cor:product}
If $\m$ is the direct product of Markov chains $\{\m_i\}$,
then  $\Gap(\m)=\min_i\Gap(\m_i).$
\end{corollary}

\subsection{One-dimensional Ising model}\label{sec:1DIsing}

As a second introductory example prior to our main application, we 
consider the one-dimensional Ising model.  Here each configuration 
$\sigma \in \Omega$ is an assignment of a ``spin'' (either +1 or -1) to 
each of $n$ vertices connected to form a line; see 
\expref{Figure}{fig:Ising}.  Let $\lambda = \eee^{-\beta}$, where $\beta > 0$ 
represents inverse temperature.  We are interested in sampling from the 
Gibbs distribution given by $\pi(\sigma) = \eee^{-\beta H(\sigma)}/Z$, 
where the Hamiltonian $H(\sigma)$ is the number of edges whose 
endpoints have different spins and $Z$ is the normalizing constant 
$\sum_{\sigma \in \Omega} \eee^{-\beta H(\sigma)}$, also called the 
\emph{partition function}. (See~\cite{LPW06} for background on the 
ferromagnetic Ising model.) 

Consider the Glauber dynamics Markov chain~$\mh$.   

\vspace{.1in}
\noindent  {\bf Glauber Dynamics $\mh$} 

\vspace{.05in}
\noindent {\tt Starting at any configuration $\sigma^0$, iterate the following:} 
\begin{itemize}
\item At time $t,$ choose a vertex $1\leq i \leq n$ uniformly at random.  
\item Set the spin of vertex $i$ to $+1$ with probability $p = (\pi(\sigma^{t, i \leftarrow +}))/(\pi(\sigma^{t, i \leftarrow +})+ \pi(\sigma^{t, i \leftarrow -}))$ where $\sigma^{t, i \leftarrow +}$ is identical to $\sigma^t$ with the spin of vertex $i$ set to $+1$ (or $-1$ for $\sigma^{t, i \leftarrow -}$).
\item Otherwise, set the spin of vertex $i$ to $-1$ with probability $1-p.$
\end{itemize}

\noindent For simplicity, we will assume that $n$ is a power of $2$. To 
apply our theorem, we decompose the state space by breaking 
configurations in half along the middle edge; again, see 
\expref{Figure}{fig:Ising}.  Transitions that fix the signs on the left are 
part of the restriction chains, and transitions that fix signs on the 
right are part of the complementary restriction chains.  Thus, our 
restrictions and complementary restrictions are both $1 \times n/2$ 
Ising models for which we can readily apply induction.  Let $a$ be the 
assignment of signs to the left $n/2$ vertices and $b$ be the signs of 
the right $n/2$ vertices.  It is straightforward to see $\sigma = 
(a,b)$ gives a unique configuration and that the state space is a 
product space.  However, $\mh$ is not a direct product of independent 
Markov chains on $a$ and $b$ because the probability of changing a sign 
of either of the middle two vertices ($n/2$ or $n/2+1$) depends on the 
sign of the other middle vertex. In order to apply 
\expref{Theorem}{thm:productspaces}, we first analyze $r(a,b) = \pi(a,b)/ 
(\pi(a)\pi(b))$ and subsequently $\epsilon$.  The techniques used here 
are similar to, but simpler than, those used \expref{Section}{sec:perm}. 

\begin{figure}[h]
  \centering
     \begin{minipage}[c]{.5\textwidth}
     \centering
    \begin{tikzpicture}[scale=.15]
      \tikzstyle{node}=[fill,circle,inner sep = 0pt, minimum size =3pt]
      \tikzstyle{tnode}=[draw=none,fill=none]
      \tikzstyle{edge}=[draw,line width = 1pt,gray]
      \node[node] (v0) at (-5,0) {};
      \node[node] (v1) at (0,0) {};
      \node[node] (v2) at (5,0) {};
      \node[node] (v3) at (10,0) {};
      \node[node] (v4) at (15,0) {};
      \node[node] (v5) at (20,0) {};
      \node[node] (v6) at (25,0) {};
      \node[node] (v7) at (30,0) {};
      
	\node[tnode] at (2.5,-3) {$a$};
	\node[tnode] at (22.5,-3) {$b$};
	\node [above] at (v0) {{\footnotesize $+$}};
	\node [above] at (v1) {{\footnotesize $+$}};
	\node [above] at (v2) {{\footnotesize $-$}};
	\node [above] at (v3) {{\footnotesize $+$}};
	\node [above] at (v4) {{\footnotesize $-$}};
	\node [above] at (v5) {{\footnotesize $+$}};	
	\node [above] at (v6) {{\footnotesize $+$}};
	\node [above] at (v7) {{\footnotesize $+$}};
        \draw[edge] (v0) -- (v1);	
        \draw[edge] (v1) -- (v2);
	\draw[edge] (v2) -- (v3);
	\draw[edge] (v3) -- (v4);
	\draw[edge] (v4) -- (v5);
	\draw[edge] (v5) -- (v6);
	\draw[edge] (v6) -- (v7);
	\draw [dashed] (12.5, 4) -- (12.5,-5);
    \end{tikzpicture}    
       \end{minipage}\hfill
       \begin{minipage}[c]{.5\textwidth}
     \centering
     {\footnotesize
     \begin{tabular}{rl}
     $w(a)$ & $= \lambda^2$\\
      $w(b)$ & $= \lambda$\\
      $c_{ab}$ & $= \lambda$\\
     \end{tabular}}
     \end{minipage}
      \caption{An example configuration of the one-dimensional Ising model.}
\label{fig:Ising}
\end{figure}


Define $\lambda = \eee^{-\beta}$. Let $w(a) = \lambda^{H(a)}$, where 
$H(a)$ is the number of edges in $a$ (the left half) with disagreeing 
signs.  Analogously, define $w(b) = \lambda^{H(b)}$.  Let $c_{ab} = 
\lambda$ if the middle signs disagree and $c_{ab} = 1$ otherwise.  Thus 
$\pi(a,b) = w(a)w(b)c_{ab}/Z$. Let $\Omega^* \subset \Omega$ be the 
configurations where the middle two vertices agree. Define $Z_{A} = 
\sum_{a}w(a)$ and $Z_B= \sum_b w(b).$  

For any fixed $a$, we have $\sum_{b: (a,b)\in\Omega^*} w(b) = \sum_{b:
(a,b)\notin\Omega^*} w(b)=\frac{1}{2}Z_B$, since we can swap spins 
on all vertices in $b$ to obtain a unique configuration $b'\in\Omega 
\setminus \Omega^*$ with $w(b')=w(b)$.  
Thus, we have
\[Z = \sum_{(a,b) \in \Omega^*}w(a) w(b) + \lambda \sum_{(a,b) \notin
    \Omega^*}w(a)w(b) = (1 + \lambda)\sum_{(a,b) \in \Omega^*}w(a)
  w(b) = \frac{(1+\lambda)Z_AZ_B}{2}\,.\]

We consider two different cases for $r(a,b)$ depending on
whether the sign of the middle two vertices agree.  First consider
the case where $(a,b) \in \Omega^*$. Here we have
\[\pi(a) = \sum_{b'}\pi(a,b') =\sum_{b': (a,b')\in
\Omega^*} \frac{w(a)w(b')}{Z} + \lambda \sum_{b': (a,b')\notin \Omega^*}
\frac{w(a)w(b')}{Z}~=~ \frac{w(a)(1+\lambda)Z_B}{2Z}\,,\]
and similarly $\pi(b) =w(b)(1+\lambda)Z_A/(2Z)$.  Therefore
\[r(a,b)= \frac{\pi(a,b) }{\pi(a)\pi(b)} = \frac{4Zw(a)w(b)
   c_{ab}}{w(a)w(b)(1+\lambda)^2Z_AZ_B} = \frac{2}{1+\lambda}\,.\]

\noindent The next case is almost identical except that $c_{ab} = \lambda$, 
so we have for $(a,b) \notin \Omega^*$ that $r(a,b) = 2\lambda/(1 + \lambda)$.

Next we use our analysis of $r(a,b)$ to bound $\epsilon$. First 
notice that since $\sum_{(a,b) \in \Omega^*} \pi(a,b) + \sum_{(a,b) 
\notin \Omega^*} \pi(a,b)  = 1$ and $\sum_{(a,b) \notin \Omega^*} 
\pi(a,b) = \lambda \sum_{(a,b) \in \Omega^*}\pi(a,b)$, we have that 
$\sum_{(a,b) \in \Omega^*} \pi(a,b) = 1/(1+\lambda)$ and $\sum_{(a,b) 
\notin \Omega^*} \pi(a,b)  = \lambda/(1 + \lambda)$. This yields

\begin{align*}
\epsilon &\leq
 \sum_{(a,b) \in \Omega^*} \pi(a,b)(\sqrt{r(a,b)} -1/\sqrt{r(a,b)})^2  + \sum_{(a,b) \notin \Omega^*} \pi(a,b)(\sqrt{r(a,b)} -1/\sqrt{r(a,b)})^2\\
&= \frac{1}{1+\lambda} \left(\sqrt{ \frac{2}{1+\lambda}} -\sqrt{ \frac{1+\lambda}{2}}\right)^2 + \frac{\lambda}{1+\lambda} \left(\sqrt{\frac{2\lambda}{1+\lambda}} -\sqrt{ \frac{1+\lambda}{2\lambda}}\right)^2 = \left(\frac{1-\lambda}{1+\lambda}\right)^2\,. \\
\end{align*}
%
Applying the complementary decomposition theorem (\expref{Theorem}{thm:productspaces}) 
gives the following recurrence: 
$\gamma_{n} \geq \gamma_{n/2}\left(\frac{2\lambda}{1+\lambda}\right)^2.$
Since the base case has gap $\Omega(n^{-1})$, this solves to $\gamma_n
= \Omega(n^{-c})$ for $c = 1 +
2\log_2\left(\frac{1+\lambda}{2\lambda}\right).$  Note that while this
does not give a tight bound, the constant $c$ is strictly better than
the constant given by~\cite{jstv04} and, unlike earlier
decomposition approaches, we have not incurred an extra factor of $n$
with each application of the decomposition theorem. 

%
\subsection{Ising model on bounded-degree trees}  %

As in~\cite{jstv04}, our proof for the one-dimensional Ising model can 
be easily generalized to trees with constant maximum degree $r$.  A 
straightforward induction shows that such a tree $T$ on $n$ vertices 
has an edge whose deletion cuts $T$ into two components, each with size 
at least $n/(r+1)$. We let $a$ represent the spins on one component and 
$b$ the spins on the other. At each level of the induction, we compute 
$r(a,b)$ and $\epsilon$ using arguments similar to those in 
\expref{Section}{sec:1DIsing} to get $\gamma_n = \Omega(n^{-c})$ for $c = 1 
+ 2\log_{(r+1)/r}\left(\frac{1+\lambda}{2\lambda}\right)$. 

\section{Generalizations and comparison with other theorems}\label{sec:generalizations}

In \expref{Section}{sec:gencomp} and~\expref{Section}{sec:classcomp}, we present 
several generalizations of \expref{Theorem}{thm:productspaces} and compare 
these results with related prior work.  Our decomposition theorems fall 
into two categories: complementary decomposition theorems that rely on 
the notion of $\epsilon$-orthogonality between the restrictions and 
complementary restrictions, and more classical decomposition theorems 
based on the projection Markov chain.  In \expref{Section}{sec:summary}, we 
summarize the proofs, with more details given in \expref{Section}{sec:decompProofs} and \expref{Section}{sec:classical}.

\subsection{Generalized complementary decomposition theorems}\label{sec:gencomp}

\expref{Theorem}{thm:productspaces} generalizes easily to non-product spaces.  
Define $r(i,j)=\pi(\Omegahat_i\cap\Omegatilde_j)/(\pi(\Omegahat_i)\pi(\Omegatilde_j))$, 
for any $1\leq i\leq r, 1\leq j\leq \rtilde$.
We say that $\{\Omega_i\}$ and $\{\Omegatilde_j\}$
is an $\epsilon$-orthogonal decomposition of $\m$ if
\begin{equation*}
  \epsilon=  \sum_{(i,j)}\pi(\Omegahat_i\cap\Omegatilde_j)(\sqrt{r(i,j)}-1/\sqrt{r(i,j)})^2\,.
\end{equation*}

\begin{theorem}\label{thm:indep}
For any $\epsilon$-orthogonal decomposition of $\m$, $\Gap(\m) 
\geq \min\{\gminhat, \gmintilde\}\left(1 -\sqrt{\epsilon}\right)^2.$
\end{theorem}

\noindent
We use \expref{Theorem}{thm:primaldual} to prove \expref{Theorem}{thm:indep},
which in turn implies \expref{Theorem}{thm:productspaces}.

\begin{theorem}\label{thm:primaldual}
$\Gap(\m)\geq \min_{x\perp \sqrt{\pi}, \|x\|=1}  
\gminhat\|\Perp{\xhat}\|^2+\gmintilde\|\Perp{\xtilde}\|^2.$
\end{theorem}

Here, $\Perp{\xhat}$ and $\Perp{\xtilde}$ are orthogonal projections of 
a vector $x$ onto the complement of the eigenspace of the top 
eigenvectors of certain matrices (defined in \expref{Section}{sec:tight}) 
containing the $P_i$'s and $\Ptilde_j$'s, respectively. This theorem is 
similar to a special case of the main result in~\cite{Dest03}. 
Destainville~\cite{Dest03} introduced a ``multi-decomposition'' scheme 
that uses $m$ different partitions of $\Omega$. In Destainville's 
result, $\|\Perp{\xhat}\|^2+\|\Perp{\xtilde}\|^2$ is replaced by a 
function of the norm of a ``multi-projection'' operator $\Pi$. Bounding 
these norms is essential, as the Markov chain $\m$ can require 
exponential time to mix even if all of the restrictions and 
complementary restrictions are polynomially mixing\footnote{Indeed, the 
introduction of the projection chain in~\cite{madr} was a key insight 
to the original decomposition theorem.}.

Unfortunately, bounding these norms can be challenging. 
Destainville~\cite{Dest03} bounds the norm of the projection $\Pi$ by 
the spectral gap of a smaller matrix $\bar{\Pi}$. In some cases, this 
gap can be analyzed directly, or even computationally for particular 
problem instances. However, for very complex distributions such as the 
distribution over biased permutations we consider here, it can be 
challenging to find the spectral gap of $\bar{\Pi}$. We believe one of 
our main contributions is the definition of $\epsilon$-orthogonality, a 
concrete combinatorial quantity that may be easier to analyze.  This 
approach is particularly useful when the chain decomposes into pieces 
that are nearly independent, as in the setting of 
\expref{Theorem}{thm:productspaces}.

\subsection{Classical decomposition theorems}\label{sec:classcomp}

The disjoint decomposition theorem of~\cite{MR00} states that the 
spectral gap $\gamma$ of $\m$ satisfies $\gamma \geq 
\frac{1}{2}\gmin\gammabar$, where, as we recall from 
\expref{Section}{sec:intro}, $\gmin=\min_i \gammahat_i$ and $\gammabar$ is 
the spectral gap of a projection chain over states $[\rhat]$. Jerrum, 
Son, Tetali, and Vigoda~\cite{jstv04} considered two quantities related 
to the spectral gap: the Poincar\'{e} and log-Sobolev constants.  
There, the authors defined a parameter 
$T=\max_i\max_{\elemone\in\Omegahat_i}\sum_{\elemtwo\in 
\Omega\setminus\Omegahat_i}\transmat(\elemone,\elemtwo)$, which can be 
seen as the maximum probability of escape from one part of the 
partition in a single step of $\transmat$, and used it to produce a 
bound on the order of the minimum gap
when $T$ is on the order of~$\gammabar$.  They also provided 
improved bounds when another parameter $\eta$ is close to zero; this 
requires a pointwise regularity condition. More recently, Pillai and 
Smith~\cite{PS17} introduced other conditions in order to directly 
bound the mixing time by a constant times the maximum of the mixing 
times of the projection and the restrictions.

The techniques developed for proving the complementary decomposition 
theorems introduced in this paper can be further applied to prove the 
following ``classical''-style decomposition theorem.

\begin{theorem}\label{thm:main}
  Let $\rho=\sqrt{2T/\gammabar}.$ Then   $\Gap(\m)\geq \displaystyle\min_{p^2+q^2=1}
\gmin q^2+\gammabar\left(q\rho - p\right)^2.$
%
\end{theorem}
\noindent
 We state a more general version of this theorem, \expref{Theorem}{thm:main2}, in
\expref{Section}{sec:classical}.
This bound allows us to rederive several known classical
decomposition theorems.

\begin{corollary}\label{cor:compareproduct}
  Assume $\m$ is lazy.  Then $\gamma\geq \gmin\gammabar/3$.
\end{corollary}
\noindent In fact, one can show that the constant is $1/2$ if $\gmin,\gammabar\leq 1/2$ (which is a common situation) or $\deltahpet\geq 1/2$ ($\deltahpet$ is defined in \expref{Section}{sec:combined}).
In \expref{Corollary}{cor:compareT} we show that \expref{Theorem}{thm:main}
can be seen as a generalization of Theorem~1 of~\cite{jstv04}, except
that it instead bounds the spectral gap.

\begin{corollary}\label{cor:compareT}
 $\gamma\geq \min\left\{\frac{\gammabar}{3}, \frac{\gmin\gammabar}{3T+\gammabar}\right\}$.
\end{corollary}

\noindent In particular, if $T/\gammabar$ is a constant, then
we get within a constant of the minimum gap as well.  \expref{Theorem}{thm:main} produces
slightly improved bounds over \expref{Corollary}{cor:compareT} when
$T\approx \gammabar\ll \gmin.$  

\subsection{Summary of the proofs of the decomposition theorems}\label{sec:summary}
  
Our proofs are elementary and use only basic facts from linear algebra 
about eigenvalues and eigenvectors.  We have chosen to assume the 
Markov chains are discrete and finite to keep the proofs as accessible 
as possible. We utilize the following standard characterization of the 
second largest eigenvalue $\lambda$ of a symmetric matrix $A$ with top 
eigenvector $v$:
\begin{equation}\label{eq:eigen} 
\lambda = \max_{x\perp v}\frac{\langle x,xA\rangle}{\|x\|^2}
           = \max_{x\perp v:\|x\|=1}\langle x,xA\rangle\,.  
\end{equation} 

\noindent For a general reversible Markov chain with transition matrix 
$\transmat$, we apply \expeqref{equation}{eq:eigen} to a symmetric matrix 
$\A=\A(\transmat)$ that has the same eigenvalues as $\transmat$. 

We apply the Vector Decomposition Method from the expander graph 
literature (see, \eg,~\cite{RVW, expander-book}), and decompose the 
vector $x$ into $\Perp{\xhat}+\Par{\xhat}$, where $\Par{\xhat}$ is 
parallel to the top eigenvector of each restriction matrix.  The 
intuition of this method is that if a particular distribution is far 
from stationarity, then it will either be far from stationarity on some 
part of the partition or on the projection, and therefore applying 
$\transmat$ brings us closer to stationarity. The benefit of this 
approach is that it allows us to quantify the independence of the 
restriction chains with the projection or complementary restriction chains.
Using 
\expeqref{equation}{eq:eigen}, for any $x\perp v$, we need to bound 
\begin{equation}\label{eq:split}
\langle x, xA\rangle = \langle \Perp{\xhat},\Perp{\xhat}A\rangle + 
\langle\Par{\xhat},\Par{\xhat}A\rangle + 2\langle \Perp{\xhat},\Par{\xhat}A\rangle\,.
\end{equation}
It is easy to bound $\langle \Perp{\xhat},\Perp{\xhat}A\rangle$ and 
$\langle\Par{\xhat},\Par{\xhat}A\rangle$ using ideas from other 
decomposition results~\cite{jstv04, MR00}. The term $\langle 
\Perp{\xhat},\Par{\xhat}A\rangle$ determines whether the decomposed 
Markov chain is nearly the direct product of two independent Markov 
chains $\m_1$ and $\m_2$, in which case $\langle 
\Perp{\xhat},\Par{\xhat}A\rangle\approx0$ and $\Gap(\m)\approx 
\min\{\Gap(\m_1),\Gap(\m_2)\}$, or whether they are far from 
independent, in which case $\langle \Perp{\xhat},\Par{\xhat}A\rangle$ 
is large and $\Gap(\m)=\Theta(\gmin \gammabar)$.  The key to our 
decomposition proofs lies in our bounds on $\langle 
\Perp{\xhat},\Par{\xhat}A\rangle$, which are different for our 
complementary decomposition theorems than they are for our classical 
decomposition theorems.   More details are provided in \expref{Section}{sec:decompProofs} and \expref{Section}{sec:classical}.

\section{Complementary decomposition theorems}\label{sec:decompProofs} 


In this section, we present the proofs of the complementary 
decomposition theorems.  First, in \expref{Section}{sec:decompPrelim}, we 
introduce some notation and terminology which will be useful for the 
proofs.  Next, we present key lemmas for both the complementary 
decomposition theorem and for the classical decomposition theorem in 
\expref{Section}{sec:combined}. In \expref{Section}{sec:orthog}, we prove our 
complementary decomposition theorems, 
\expref{Theorems}{thm:productspaces},~\expref{}{thm:indep}, 
\expref{and}{thm:primaldual}.  

%
\subsection{Extended preliminaries}  %
\label{sec:decompPrelim}

We first fix some notation and terminology. The symbol $\otimes$ is 
used for tensor product.  We write $(v)_i$ to mean the $i^{\text{th}}$ 
coordinate of a vector $v$. The second largest eigenvalue of $P_i$ will 
be denoted $\lambda_i$, and $\lambda_{\max} = \max_i \lambda_i$. The 
``top eigenvector'' of a matrix will be the eigenvector corresponding 
to the eigenvalue of largest absolute value.

In order to prove our decomposition results, we wish to apply 
\expeqref{equation}{eq:eigen} to $\transmat$.  However, since $\transmat$ may 
not be symmetric, we appeal to the following symmetrization technique 
that appears in~\cite[p. 153]{LPW06}. Given $\transmat$ with stationary 
distribution $\pi$, define a matrix $\A := \A(\transmat)$ by 
$\A(\elemone,\elemtwo) :=\pi(\elemone)^{1/2}\pi(\elemtwo)^{-1/2}\transmat(\elemone,\elemtwo)$. 
$\A$ is similar to $\transmat$ (\ie,  they have the same eigenvalues), 
but is symmetric, so we can infer a bound on the second eigenvalue of 
$\transmat$ by applying \expeqref{equation}{eq:eigen} to $\A$. It is easy to 
check that the top eigenvector of $\A$ is $\sqrt{\pi}$, which is the 
vector with entries $\sqrt{\pi(\elemone)}$ for any $\elemone\in\Omega$. 

We apply this same symmetrization technique to other matrices as well. 
For $i \in [\rhat]$ we let $A_i := \A(P_i)$ and for $i \in 
[\rtilde]$ we let $\Atilde_i := \A(\Ptilde_i)$.  We then write $\Ahat$ 
to mean the $|\Omega| \times |\Omega|$ matrix with 
$\Ahat(\elemone,\elemtwo) = A_i(\elemone,\elemtwo)$ if 
$\elemone,\elemtwo \in \Omegahat_i$ for some $i \in [\rhat]$, and zero 
otherwise. Analogously, we write $\Atilde$ to mean the $|\Omega| \times 
|\Omega|$ matrix with $\Atilde(\elemone,\elemtwo) = 
\Atilde_i(\elemone,\elemtwo)$ if $\elemone,\elemtwo \in \Omegatilde_i$ 
for some $i \in [\rtilde]$, and zero otherwise. It is important to note 
that $\Ahat \not= \A(\Phat)$ and 
$\Atilde \not= \A(\Ptilde)$.  
\begin{proposition}\label{prop:matdecomp}
The matrix $\A$ satisfies 
$\A = \Ahat+\Atilde-I_{|\Omega|}.$
\end{proposition}

Let $\lambdatilde_1\geq\lambdatilde_2\geq\ldots\geq 
\lambdatilde_{|\Omega|}$ be the eigenvalues of $\Atilde$ with 
corresponding eigenvectors $\vtilde_1,\vtilde_2,\ldots, 
\vtilde_{|\Omega|}$.  As $\Atilde$ is symmetric, the real spectral 
theorem tells us that its eigenvectors form an orthonormal basis of 
$\R^{|\Omega|}$.  We consider the basis representations $\Perp{\xhat} = 
\sum_i\Perp{a_i}\vtilde_i$ and $\Par{\xhat} = 
\sum_i\Par{a_i}\vtilde_i$.  More generally, for any $v \in 
\R^{|\Omega|}$, we write  $v=\sum_{i}a_i\vtilde_i$ for some constants 
$a_1,a_2,\ldots, a_{|\Omega|} \in \R$.  Also, $\|v\|^2 = \sum_i 
a_i^2\|\vtilde_i\|^2,$ and $v\Atilde = \sum_i 
a_i\lambdatilde_i\vtilde_i$.

\subsection{Key ideas and lemmas for the proofs}\label{sec:combined}

We wish to apply \expeqref{equation}{eq:eigen} to $\A$. Recall that 
$\sqrt{\pi}$ is the top eigenvector of $\A$.  Let $x \in \R^{|\Omega|}$ 
with $x\perp \sqrt{\pi}$ and $\|x\|=1$. We will decompose $x$ into two 
vectors $\Par{\xhat}$ and $\Perp{\xhat}$ as follows (note: this is 
similar to the vector decomposition used for the Zig Zag Product 
in~\cite{RVW}). For any $i\in [\rhat]$, let 
$\xhat_i\in\R^{|\Omegahat_i|}$ be the vector defined by 
$\xhat_i(\elemone)=\xhat(\elemone)$ for all $\elemone \in \Omegahat_i$. 
Then $x=\sum_i e_i\otimes \xhat_i$.  We further decompose $\xhat_i$ 
into $\Par{\xhat_i}$, the part that is parallel to $\sqrt{\pihat_i}$, 
and $\Perp{\xhat_i}$, the part that is perpendicular to 
$\sqrt{\pihat_i}$; recall that $\sqrt{\pihat_i}$ is the top eigenvector 
of $A_i$. Finally, define $\Par{\xhat}, \Perp{\xhat} \in \R^{|\Omega|}$ 
by $\Par{\xhat}=\sum_i e_i\otimes \Par{\xhat_i}$ and 
$\Perp{\xhat}=\sum_i e_i\otimes\Perp{\xhat_i}$.  Hence $x = \sum_i 
e_i\otimes \xhat_i = \Par{\xhat}+\Perp{\xhat}$. Define $\Par{\xtilde}$ 
and $\Perp{\xtilde}$ analogously.

As described in \expref{Section}{sec:summary}, we will bound $\langle 
x,xA\rangle$ via \expeqref{equation}{eq:split}: \[\langle x, xA\rangle = 
\langle \Perp{\xhat},\Perp{\xhat}A\rangle + 
\langle\Par{\xhat},\Par{\xhat}A\rangle + 2\langle 
\Perp{\xhat},\Par{\xhat}A\rangle\,.\] We need the following simple 
proposition.

\begin{lemma}\label{lem:xparA}  
The following holds: $\Par{\xhat}\A=\Par{\xhat}\Atilde.$
\end{lemma}

\noindent Applying \expref{Lemma}{lem:xparA}, \expeqref{equation}{eq:split} becomes
\[ \langle {x},{x}A\rangle=  \langle \Par{\xhat},\Par{\xhat}\Atilde\rangle +2
  \langle\Perp{\xhat},\Par{\xhat}\Atilde\rangle
  +\langle\Perp{\xhat},\Perp{\xhat}(\Ahat+\Atilde-I_{|\Omega|})\rangle\,.\]

For ease of notation, we define the following quantities:
\[ \deltahpeh = \frac{\langle\Perp{\xhat},\Perp{\xhat}\Ahat\rangle}{\|\Perp{\xhat}\|^2},~ 
\deltahpet = \frac{\langle\Perp{\xhat},\Perp{\xhat}\Atilde\rangle}{\|\Perp{\xhat}\|^2},~
\deltahpah = \frac{\langle\Par{\xhat},\Par{\xhat}\Ahat\rangle}{\|\Par{\xhat}\|^2},~
\deltahpat =
\frac{\langle\Par{\xhat},\Par{\xhat}\Atilde\rangle}{\|\Par{\xhat}\|^2}. \]
Plugging these in, we have
\begin{align}
 \langle {x},{x}A\rangle&=  \deltahpat \|\Par{\xhat}\|^2 +2 \langle\Perp{\xhat},\Par{\xhat}\Atilde\rangle
  +(\deltahpeh +\deltahpet -1)\|\Perp{\xhat}\|^2\,.%
\end{align}

Bounding  $\deltahpeh$ and $\deltahpat$ is straightforward, and borrows 
many of the ideas from classical decomposition results.  If 
$\Par{\xhat}\Atilde$ were orthogonal to $\Perp{\xhat}$, then doing so 
would be sufficient to proving a strong decomposition theorem.  
However, this is not true in general, so we must also bound 
$\langle\Perp{\xhat},\Par{\xhat}\Atilde\rangle$.  Our two types of 
theorems do so in different ways, which are presented in 
\expref{Section}{sec:orthog} and \expref{Section}{sec:classical}.

The next lemma makes concrete the intuition that if a particular 
distribution is far from stationarity, then it will either be far from 
stationarity on some restriction---in which case $\Ahat$ will bring it 
closer to stationarity (as in part 1)---or on the projection---in which 
case $\Atilde$ will bring it closer to stationarity (as in part 2).  
The proof is straightforward from the definitions. 

\begin{lemma}\label{lem:eta}
  With the above notation,
\begin{enumerate}
\item $\deltahpeh \leq \lambda_{\max}.$
\item $\deltahpat \leq \lambdabar.$
\end{enumerate}
\end{lemma}
\noindent Note that $\lambda$ is formally defined in \expeqref{equation}{eq:eigen} and $\lambda_{\max}$ is defined at the beginning of \expref{Section}{sec:decompPrelim}.

\subsection{Proofs of complementary decomposition theorems}\label{sec:tight}\label{sec:orthog}

Next we use the technology developed in \expref{Section}{sec:combined} to 
prove \expref{Theorem}{thm:primaldual}. Recall that 
$\lambdatilde_1\geq\lambdatilde_2\geq\ldots\geq 
\lambdatilde_{|\Omega|}$ are the eigenvalues of $\Atilde$ with 
corresponding eigenvectors $\vtilde_1,\vtilde_2,\ldots, 
\vtilde_{|\Omega|}$, and that for any $v \in \R^{|\Omega|}$, we write 
$v=\sum_{i}a_i\vtilde_i$ for some constants $a_1,a_2,\ldots, 
a_{|\Omega|} \in \R$.  Also, $\|v\|^2 = \sum_i a_i^2\|\vtilde_i\|^2,$ 
and $v\Atilde = \sum_i a_i\lambdatilde_i\vtilde_i$. Define the set 
$S=\{i:\mu_i=1\}$. Let 
$\deltatpet = \frac{\langle\Perp{\xtilde},\Perp{\xtilde}\Atilde\rangle}{\|\Perp{\xtilde}\|^2}$.
 Now we can make an explicit statement about the gap 
of $\m$; notice the equality in \expeqref{equation}{eq:primaldual}.  
\begin{customthm}{\ref{thm:primaldual}}
\begin{equation}\label{eq:primaldual}
\Gap(\m)= \min_{x\perp \sqrt{\pi}, \|x\|=1} (1-\deltahpeh)\|\Perp{\xhat}\|^2+
(1-\deltatpet)\|\Perp{\xtilde}\|^2\,.
\end{equation}
In particular,
\begin{equation}\label{eq:primaldual2}
\Gap(\m)\geq \min_{x\perp \sqrt{\pi}, \|x\|=1}  
\gminhat\|\Perp{\xhat}\|^2+\gmintilde\|\Perp{\xtilde}\|^2\,.
\end{equation}
\end{customthm}

\begin{proof}
Notice $\deltatpet\|\Perp{\xtilde}\|^2 = \sum_{i\in S}\lambdatilde_i(\Perp{a_i}+\Par{a_i})^2$.
Thus,
\[(1-\deltatpet)\|\Perp{\xtilde}\|^2 = \sum_{i}(1-\lambdatilde_i)(\Perp{a_i}+\Par{a_i})^2
   = (1-\deltahpet)\|\Perp{\xhat}\|^2 + (1-\deltahpat)\|\Par{\xhat}\|^2 - 
   2\langle \Perp{\xhat},\Par{\xhat}\Atilde\rangle\,.\]
On the other hand, from \expeqref{equation}{eq:strongestbound}, we have
\[1-\langle x,x\A\rangle = (1-\deltahpeh)\|\Perp{\xhat}\|^2+(1-\deltahpet)\|\Perp{\xhat}\|^2+
(1-\deltahpat)\|\Par{\xhat}\|^2 - 2\langle \Perp{\xhat},\Par{\xhat}\Atilde\rangle\,.\]
Thus, for all $x\perp \sqrt{\pi}$ with norm 1, we have
\[1-\langle x,x\A\rangle=(1-\deltahpeh)\|\Perp{\xhat}\|^2+(1-\deltatpet)\|\Perp{\xtilde}\|^2\,.\]
Applying \expeqref{equation}{eq:eigen}, we get \expeqref{equation}{eq:primaldual}.
To get \expeqref{equation}{eq:primaldual2}, we apply \expref{Lemma}{lem:eta}, which yields
$1-\deltahpeh \geq 1-\lambda_{\max} = \gminhat$.  An analogous statement
to \expref{Lemma}{lem:eta} holds for $\deltatpet$, and shows 
$1-\deltatpet \geq \gmintilde$. 
\end{proof}

It remains to prove \expref{Theorem}{thm:indep}.
By \expref{Theorem}{thm:primaldual}, if $\gminhat$ and $\gmintilde$ are not 
too small, it suffices to show that $\|\Perp{\xhat}\|^2$ and 
$\|\Perp{\xtilde}\|^2$ cannot both be small.  To this end, we further 
decompose $\Perp{\xhat}$ and $\Par{\xtilde}$ based on the eigenvectors 
of $\Atilde$.  Define $S=\{i:\mu_i=1\}$ and vectors $\b=\sum_{i\in S} \Par{a_i}v_i$ and 
$\a=\sum_{i\notin S} \Par{a_i}v_i$. Similarly, let $\d=\sum_{i\in S} 
\Perp{a_i}v_i$ and $\c=\sum_{i\notin S} \Perp{a_i}v_i$. Notice 
$\Par{\xtilde} = \b+\d$ and $\Perp{\xtilde} =\a+\c$, so that the 
vectors in each row (respectively, column) of the following table sum to the 
vector in its row (respectively, column) label.
\begin{center}
\begin{tabular}{c|c|c|}
\multicolumn{1}{c}{}
 & \multicolumn{1}{c}{$\Par{\xtilde}$}
 & \multicolumn{1}{c}{$\Perp{\xtilde}$} \\
\cline{2-3}
$\Par{\xhat}$ &  $\b$ & $\a$\\
\cline{2-3}
$\Perp{\xhat}$ & $\d$ & $\c$\\
\cline{2-3}
\end{tabular}
\end{center}
The vectors within each row are orthogonal, as they are in the span of 
eigenvectors with distinct eigenvalues.  However, the vectors within 
each column are not necessarily orthogonal.

The idea of the proof of \expref{Theorem}{thm:indep} is that if
$\|\b\|$ is small, then $\|\Perp{\xhat}\|^2+ \|\Perp{\xtilde}\|^2$ is
large.    The following lemma states that 
$\epsilon$-orthogonality is sufficient to guarantee $\|\b\|$ is small.

\begin{lemma}\label{lem:directproduct}
Let $\epsilon$ be as defined in \expeqref{equation}{eq:eps}.
Then 
  $\displaystyle \|\b\|^2\leq \epsilon.$
\end{lemma}
\begin{proof}
Recall $\b$ is the projection of $\Par{\xhat}$ onto the top
eigenvectors of $\Atilde$.  The top eigenvectors of $\Atilde$ are precisely
the set of all $\sqrt{\pitilde_j}$ for $j\in[\rtilde]$.
Therefore,
\[\b = \sum_{j}\frac{\langle
\Par{\xhat},\sqrt{\pitilde_j}\rangle}{\|\sqrt{\pitilde_j}\|^2}\sqrt{\pitilde_j}\,.\]
As the eigenvectors of $\Atilde$ are an orthonormal basis, we have
\[\|\b\|^2 = \sum_{j}\langle \Par{\xhat},\sqrt{\pitilde_j}\rangle^2\,.\]
For any $j\neq j'\in[\rtilde]$ and any $\sigma\in\Omegatilde_{j'}$, $\sqrt{\pitilde_j(\sigma)}=0$  and for $i\in[\rhat]$, $\pitilde_j(\Omegahat_i\cap\Omegatilde_j) = \pi(\Omegahat_i\cap\Omegatilde_j)/\pi(\Omegatilde_j)$.  Therefore,
\begin{equation}\label{eq:innerprod}
 \langle \Par{\xhat},\sqrt{\pitilde_j}\rangle 
 ~= \sum_i\sum_{\sigma\in\Omegahat_i}\Par{\xhat}(\sigma)\sqrt{\pitilde_j(\sigma)}
 ~= \sum_i\alpha_i\sum_{\sigma\in\Omegahat_i\cap\Omegatilde_j}\sqrt{\pihat_i(\sigma)\pitilde_j(\sigma)}
 ~= \sum_{i\in[\rhat]}\alpha_i\frac{\pi(\Omegahat_i\cap\Omegatilde_j)}{\sqrt{\pi(\Omegahat_i)\pi(\Omegatilde_j)}}\,.
 \end{equation}

Since $x \perp \sqrt{\pi}$ and $\Perp{\xhat} \perp \sqrt{\pi}$
by definition, it follows that $\Par{\xhat} \perp \sqrt{\pi}$ as well.
This implies that $\alpha \perp \sqrt{\pibar}$, as 
\begin{equation}\label{eq:alphaperp}
0 = \langle \Par{\xhat},\sqrt{\pi} \rangle
  = \sum_i \alpha_i \sum_{\elemone \in \Omegahat_i} \sqrt{\pihat_i(\elemone)\pi(\elemone)}
  = \sum_i \alpha_i \sum_{\elemone \in \Omegahat_i} \frac{\sqrt{\pi(\elemone)}}{\sqrt{\pi(\Omegahat_i)}} \sqrt{\pi(\elemone)} 
  = \sum_i \alpha_i \sqrt{\pi(\Omegahat_i)}\,,
\end{equation}
and this final term is equal to $\sum_i \alpha_i \sqrt{\pibar_i}=\langle \alpha, 
\sqrt{\pibar} \rangle$.  Multiplying \expeqref{equation}{eq:alphaperp} by 
$\pi(\Omegatilde_j)$ and subtracting it from \expeqref{equation}{eq:innerprod}, we have
\[
 \langle \Par{\xhat},\sqrt{\pitilde_j}\rangle 
 =\sum_{i\in[\rhat]}\alpha_i\left(\frac{\pi(\Omegahat_i\cap\Omegatilde_j)}{\sqrt{\pi(\Omegahat_i)\pi(\Omegatilde_j)}}-
    \sqrt{\pi(\Omegahat_i)\pi(\Omegatilde_j)}\right)
 =\langle\alpha,V_j\rangle,
 \]
where 
\[V_j(i):=\left(\frac{\pi(\Omegahat_i\cap\Omegatilde_j)}{\sqrt{\pi(\Omegahat_i)\pi(\Omegatilde_j)}}-
\sqrt{\pi(\Omegahat_i)\pi(\Omegatilde_j)}\right)=\sqrt{\pi(\Omegahat_i\cap\Omegatilde_j)}(\sqrt{r(i,j)}-1/\sqrt{r(i,j)})\,.\]  
By the Cauchy--Schwarz inequality and
the fact that $\|\alpha\|=\|\Par{\xhat}\| \leq \|x\|=1$, we have
$ \langle\alpha,V_j\rangle\leq \|\alpha\|\|V_j\| = \|V_j\|.$
Therefore we get,
\[ \|\b\|^2 = \sum_{j}\langle \Par{\xhat},\sqrt{\pitilde_j}\rangle^2
	\leq \sum_j\|V_j\|^2 
	= \sum_{i,j}\pi(\Omegahat_i\cap\Omegatilde_j)(\sqrt{r(i,j)}-1/\sqrt{r(i,j)})^2\,.\qedhere\]
\end{proof}

To prove \expref{Theorem}{thm:indep} from \expref{Theorem}{thm:primaldual}, we 
must show that if $\|\b\|^2\leq \epsilon$, then $\|\Perp{\xhat}\|^2+ 
\|\Perp{\xtilde}\|^2\geq (1-\sqrt{\epsilon})^2$.  As the sum of the 
squared norms of the vectors in the above table is 1, it is reasonable 
to expect that if $\|\b\|^2$ is small, then $\|\Perp{\xhat}\|^2+ 
\|\Perp{\xtilde}\|^2$ is large.  However, this is not as 
straightforward as one might expect, as the vectors within each column 
are not necessarily orthogonal, so we may have $\|\Perp{\xtilde}\|^2< 
\|\a\|^2 + \|\c\|^2$. 

%

%
%
\begin{proof}[Proof of \expref{Theorem}{thm:indep}]   %
By \expeqref{equation}{eq:primaldual2} from \expref{Theorem}{thm:primaldual}, 
it suffices to show that $\|\Perp{\xhat}\|^2+ \|\Perp{\xtilde}\|^2\geq
(1-\sqrt{\epsilon})^2$. If $\|\Perp{\xhat}\|^2\geq (1-\sqrt{\epsilon})^2$ we are done,
so we may assume not.
As the vectors within each row of the table are orthogonal, we have
$(1-\sqrt{\epsilon})^2>\|\Perp{\xhat}\|^2 =\|\c\|^2+ \|\d\|^2$.
Furthermore, since
\[0=\langle\Perp{\xhat}, \Par{\xhat}\rangle =\sum_i
\Perp{a_i}\Par{a_i}\|v_i\|^2\]
we have
\[\langle\c,\a\rangle = \sum_{i\notin S}\Perp{a_i}\Par{a_i}\|v_i\|^2 = 
-\sum_{i \in S}\Perp{a_i}\Par{a_i}\|v_i\|^2 = -\langle\d,\b\rangle\,.\]
Thus,
\begin{align}
    \|\Perp{\xtilde}\|^2 &=\|\c+\a\|^2\nonumber\\
    %
    &=\|\c\|^2+\|\a\|^2+2\langle\c,\a\rangle\label{eq:innerabcd}\\
    &=\|\c\|^2+\|\a\|^2-2\langle\d,\b\rangle\nonumber\\
    &\geq\|\c\|^2+\|\a\|^2-2\|\d\|\|\b\|\,.\nonumber
\end{align}
We used Cauchy--Schwarz for the final inequality.  Putting everything 
together,
\begin{align*}
\|\Perp{\xhat}\|^2+ \|\Perp{\xtilde}\|^2 &\geq\|\c\|^2+ \|\d\|^2 +\|\c\|^2+\|\a\|^2-2\|\d\|\|\b\|\\
  &=\|\c\|^2+ \|\d\|^2 + (1-\|\d\|^2-\|\b\|^2)-2\|\d\|\|\b\|\\
  &= 1 + \|\c\|^2-\|\b\|^2-2\|\d\|\|\b\|\\
  &\geq 1 + \|\c\|^2-\epsilon-2\sqrt{\epsilon}\sqrt{(1-\sqrt{\epsilon})^2-\|\c\|^2}\\
  &\geq 1 -\epsilon-2\sqrt{\epsilon}(1-\sqrt{\epsilon})\\
  &= (1 -\sqrt{\epsilon})^2.\qedhere
\end{align*}
\end{proof}
%

\section{Classical decomposition theorems}\label{sec:classical}
In this section, we will prove our classical decomposition theorem,
\expref{Theorem}{thm:main}.  It will be a corollary of the following theorem. 

\begin{theorem}\label{thm:main2}
  Let $\rho=\sqrt{(1-\deltahpet)/\gammabar}.$ Then  %
%
\[\gamma(\m) \geq \min_{p^2+q^2=1}\gmin q^2+\left(q\sqrt{1-\deltahpet} - p\sqrt{\gammabar}\right)^2\,.\]
\end{theorem}

\noindent With the technology developed in \expref{Section}{sec:combined}, there is
one critical piece remaining to prove \expref{Theorem}{thm:main2}, which
is to bound the cross terms generated by applying the
matrix $\Atilde$ to $\Par{\xhat}$.

\begin{lemma}\label{lem:mixed}
With the above notation, 
$|\langle\Perp{\xhat},\Par{\xhat}\Atilde\rangle| \leq 
\sqrt{(1-\deltahpat)(1-\deltahpet)}\|\Par{\xhat}\|\|\Perp{\xhat}\|.$
\end{lemma}

\begin{proof}
 Recall that $\{\lambdatilde_i\}$ are the eigenvalues
of $\Atilde$ with corresponding eigenvectors $\{\vtilde_i\}$, and
$\Perp{\xhat} = \sum_i\Perp{a_i}\vtilde_i$ and $\Par{\xhat} =
\sum_i\Par{a_i}\vtilde_i$ are the basis representations of
$\Perp{\xhat}$ and $\Par{\xhat}.$

 Since $\Perp{\xhat}$ is 
perpendicular to $\Par{\xhat}$, we have 
$0=\langle\Perp{\xhat},\Par{\xhat}\rangle = 
\sum_{i}\Perp{a_i}\Par{a_i}\|\vtilde_i\|^2$. Notice
%
%
%
%
\begin{align*}
	\langle\Perp{\xhat},\Par{\xhat}\Atilde\rangle 
	&= \sum_{i}\lambdatilde_i\Perp{a_i}\Par{a_i}\|\vtilde_i\|^2 \\
	&	= \sum_{i}\Perp{a_i}\Par{a_i}\|\vtilde_i\|^2 - \sum_{i}(1-\lambdatilde_i)\Perp{a_i}\Par{a_i}\|\vtilde_i\|^2 \\
	&	= 0 - \sum_{i}(1-\lambdatilde_i)\Perp{a_i}\Par{a_i}\|\vtilde_i\|^2\,. 
\end{align*}   %
Define vectors $\Par{\bigx}:=\sum_i \sqrt{1-\lambdatilde_i}\Par{a_i}\vtilde_i$ and
$\Perp{\bigx}:= \sum_i \sqrt{1-\lambdatilde_i}\Perp{a_i}\vtilde_i$.   Then because the
$\vtilde_i$ are mutually orthogonal, we have
\[ \langle \Perp{\bigx},\Par{\bigx}\rangle = \sum_i (1-\lambdatilde_i)\Perp{a_i}\Par{a_i}\|\vtilde_i\|^2 
	= - \langle \Perp{\xhat},\Par{\xhat}\A \rangle\,.\]
By Cauchy--Schwarz, $|\langle \Perp{\bigx},\Par{\bigx}\rangle|\leq
\|\Perp{\bigx}\|\|\Par{\bigx}\|$.  Moreover,
\[ \|\Par{\bigx}\|^2 = \sum_i (1-\lambdatilde_i)(\Par{a_i})^2\|\vtilde_i\|^2 
	= \|\Par{\xhat}\|^2 - \sum_i \lambdatilde_i(\Par{a_i})^2\|\vtilde_i\|^2
	= \|\Par{\xhat}\|^2 - \langle \Par{\xhat}, \Par{\xhat}\Atilde \rangle
	= (1-\deltahpat)\|\Par{\xhat}\|^2. \]
Similarly, $\|\Perp{\bigx}\|^2 = (1-\deltahpet)\|\Perp{\xhat}\|^2$.  Taken
together, this proves the lemma.
\end{proof}

Finally, we are ready to prove \expref{Theorem}{thm:main2}. 

%
%
%
\begin{proof}[Proof of \expref{Theorem}{thm:main2}]
Let $x\in \R^{|\Omega|}$ with $x\perp \sqrt{\pi}$ and $\|x\|=1$.  By \expeqref{equation}{eq:eigen}, 
$\Gap(\m) \geq 1-\langle x,x\A\rangle$.  
Applying \expref{Lemma}{lem:xparA} and the definitions of $\deltahpeh$, $\deltahpat$,
and $\deltahpet$, \expeqref{equation}{eq:split} becomes
\begin{align}
 \langle {x},{x}A\rangle&=  \langle \Par{\xhat},\Par{\xhat}\Atilde\rangle +2
  \langle\Perp{\xhat},\Par{\xhat}\Atilde\rangle
  +\langle\Perp{\xhat},\Perp{\xhat}(\Ahat+\Atilde-I_{|\Omega|})\rangle\nonumber\\
  &= \deltahpat \|\Par{\xhat}\|^2 +2 \langle\Perp{\xhat},\Par{\xhat}\Atilde\rangle
  +(\deltahpeh +\deltahpet -1)\|\Perp{\xhat}\|^2\,.\label{eq:strongestbound}
\end{align}
Applying \expref{Lemma}{lem:mixed}, we have
\[ \Gap(\m) \geq 1 - (\deltahpat \|\Par{\xhat}\|^2 + 
2 \sqrt{(1-\deltahpat)(1-\deltahpet)}\|\Par{\xhat}\|\|\Perp{\xhat}\|
+(\deltahpeh +\deltahpet -1)\|\Perp{\xhat}\|^2)\,.
\]
Rearranging terms and using the fact that
$1=\|x\|^2=\|\Perp{x}\|^2+\|\Par{x}\|^2$, we have
\begin{equation}\label{eq:main}
\Gap(\m) \geq \min_{x\perp\sqrt{\pi}:\|x\|=1} 
(1-\deltahpeh)\|\Perp{\xhat}\|^2+\left(\sqrt{1-\deltahpet}\|\Perp{\xhat}\| - \sqrt{1-\deltahpat}\|\Par{\xhat}\|\right)^2\,.
\end{equation}

By setting $q=\|\Perp{\xhat}\|$ and $p=\|\Par{\xhat}\|$, we immediately get
\[\Gap(\m) \geq \min_{p^2+q^2=1} 
(1-\deltahpeh)q^2+\left(\sqrt{1-\deltahpet}q - \sqrt{1-\deltahpat}p\right)^2\,.\] 


By \expref{Lemma}{lem:eta} and the definition of the spectral gap, $(1-\deltahpeh)\geq \gmin$ and
 $(1-\deltahpat)\geq \gammabar$.   Thus we have, for any values of $q$ and $p$, 
\begin{align*}
(1-\deltahpeh)q^2+\left(\sqrt{1-\deltahpet}q - \sqrt{1-\deltahpat}p\right)^2
&\geq \gmin q^2+\left(\sqrt{1-\deltahpet}q - \sqrt{1-\deltahpat}p\right)^2.
\end{align*}


Define $f(q,p) = a q^2 + (bq-cp)^2$, where $a,b,c\geq 0$.  Assume $c\geq c_*\geq 0$, and define $f'= a q^2 + (bq-c_*p)^2$.  We want to show that
\begin{equation}\label{eq:lowerbound}
 \min_{q^2+p^2=1} f \geq  \min_{q^2+p^2=1} f'\,.
\end{equation}

If $c_*=0$ then $f'$ is minimized at the point $(0,1)$, and $ f'(0,1)=0\leq  \min_{q^2+p^2=1} f$.  Thus, we may assume $c\geq c_*>0$.  If $a=0$ then $\min_{q^2+p^2=1} f =  \min_{q^2+p^2=1}(bq-cp)^2 = 0$, since we can choose $p/q = b/c$.  Similarly, in this case, $\min_{q^2+p^2=1} f'=0$.  Hence we may assume $a>0$.


The minimum of the function $f$ with respect to $q$ and $p=\sqrt{1-q^2}$ occurs either at the endpoints $q=0$ or $q=1$, or at the critical point $(q_*, p_*)$ satisfying $\frac{\partial f}{\partial q}(q_*,p_*)=0$.
Note that $f(0,1) = c^2\geq c_*^2=f'(0,1)\geq \min_{q^2+p^2=1} f'$, and $f(1,0) = a+b^2=f'(1,0)\geq \min_{q^2+p^2=1} f'$.  
Now we may assume $q, p>0$.   

Now,
\[\frac{\partial f}{\partial q}  = 2aq + 2(bq-cp)\left(b+\frac{cq}{p}\right)\,.\]
Since $c,q,p>0$ and $b\geq 0$, we have $b+\frac{cq}{p}>0$.  Thus, the point $(q_*,p_*)$  satisfies
\[bq_*-cp_* = \frac{-aq_*}{b+\frac{cq_*}{p_*}}< 0\,,\]
since $a,q_*>0$.  



%
We will now take the derivative of $f$ with respect to $c$:
\[\frac{\partial f}{\partial c}  = 2(cp-bq)p\,.\]
Thus, $\frac{\partial f}{\partial c}(q_*,p_*) = 2(cp_*-bq_*)p_* >0$.  This implies $f$ is increasing with $c$ in the neighborhood around the point $(q_*, p_*)$, and so by decreasing $c$ by $\epsilon$ we will lower the value of $\min_{q^2+p^2=1} f$; this is true for every $c>0$, so we have $\min_{q^2+p^2=1} f \geq \min_{q^2+p^2=1} f'$.  This proves \expref{Equation}{eq:lowerbound}.
%
%


Therefore, we have
\[\gamma(\m) \geq  \min_{p^2+q^2=1}\gmin q^2+\left(q\sqrt{1-\deltahpet} - p\sqrt{\gammabar}\right)^2\,.\qedhere\]
\end{proof}
%

The statement of \expref{Theorem}{thm:main2} is admittedly technical.
However, from it we may derive several corollaries, as listed in \expref{Section}{sec:generalizations}. 

To see that \expref{Theorem}{thm:main} follows from
\expref{Theorem}{thm:main2}, we prove that $\deltahpet\geq 1-2T$.

%
%
%
\begin{proof}[Proof of \expref{Theorem}{thm:main}]
Recall $T = 
\max_i\max_{\sigma\in\Omegahat_i}\sum_{\tau\in\Omega\setminus\Omegahat_i}\transmat(\sigma,\tau)$.  
This is the parameter $\gamma$ from~\cite{jstv04}. We will now show that $\deltahpet\geq 1-2T$.
Notice the probability of a move in $\Ahat$ is at least 
$1-T$, so every element has a self-loop probability of at least $1-T$ 
in $\Atilde$. 
Thus, $\Atilde$ can be 
written as $\Atilde = T\Atilde' + I(1-T)$ for some transition matrix 
$\Atilde'$ with minimum eigenvalue $-1$.  This implies that the 
minimum eigenvalue of $\Atilde$ satisfies $\mu_{\min}\geq 1-2T$.  On 
the other hand, $\deltahpet\geq \mu_{\min}$. 
\end{proof}
%

Next we prove \expref{Corollary}{cor:compareproduct}, which states that $\gamma\geq \gmin\gammabar/3$.
In fact, one can show that the constant is $1/2$ if $\gmin,\gammabar\leq 1/2$ or
$1-\deltahpet\leq 1/2$.

%
%
%
\begin{proof}[Proof of \expref{Corollary}{cor:compareproduct}]
Since the bound in \expref{Theorem}{thm:main} is minimized when 
$\deltahpet$ is minimized, we may assume $\deltahpet=0$.
Thus, $\rho=1/\sqrt{\gammabar}$.  We will show that for all $p,q$
satisfying $p^2+q^2=1$, we have
\[\gmin q^2+\gammabar\left(p - \frac{q}{\sqrt{\gammabar}}\right)^2\geq
\frac{ \gmin\gammabar}{3}\,.\]
Clearly if $q^2\geq \gammabar/3$, then we are done.  So we may assume
$q^2/\gammabar<1/3$.  Notice that since $\gammabar\leq 1$,
\[\left(p - \frac{q}{\sqrt{\gammabar}}\right)^2
~=~ \left(\sqrt{1-q^2} - \frac{q}{\sqrt{\gammabar}}\right)^2
~\geq~ \left(\sqrt{1-\frac{q^2}{\gammabar}} -
\frac{q}{\sqrt{\gammabar}}\right)^2\,.\]
As $q^2/\gammabar<1/3$, we have
$\left(\sqrt{1-\frac{q^2}{\gammabar}} -
\frac{q}{\sqrt{\gammabar}}\right)^2\geq \frac{1}{3} -
\frac{q^2}{\gammabar}.$  Therefore,
\begin{align*}
  \gmin q^2+\gammabar\left(p -  \frac{q}{\sqrt{\gammabar}}\right)^2
  &\geq  \gmin q^2+\frac{\gammabar}{3} -q^2\\
  &= \frac{\gammabar}{3} - q^2(1-\gmin)\\
  &> \frac{\gammabar}{3}(1-(1-\gmin)) \\
  &= \gmin\gammabar/3\,.\qedhere
\end{align*}
\end{proof}
%

Next, we will show that a variant of Theorem~1 of~\cite{jstv04} follows from
\expref{Theorem}{thm:main} from this paper, which is the content of
\expref{Corollary}{cor:compareT}.  \expref{Theorem}{thm:main} produces
slightly improved bounds over \expref{Corollary}{cor:compareT} when
$T\approx \gammabar\ll \gmin$.  

%
%
%
\begin{proof}[Proof of \expref{Corollary}{cor:compareT}]
We wish to show $\gamma\geq \min\left\{\frac{\gammabar}{3}, 
\frac{\gmin\gammabar}{3T+\gammabar}\right\}$. As $\rho^2\leq 
2T/\gammabar$, it suffices to show that for all $p^2+q^2=1$,
 \[\gmin q^2+\gammabar(p-\rho q)^2\geq \min\left\{\frac{\gammabar}{3}, \frac{\gmin}{1+1.5\rho^2}\right\}\,.\]
 If $q^2\geq \frac{1}{1+1.5\rho^2}$ then we are done, so we may assume
 $\frac{1}{1+1.5\rho^2}-q^2\geq 0$.
 Define
  \[f_1=\frac{(p-\rho q)^2}{1/(1+1.5\rho^2) - q^2} \quad \text{and}
  \quad f_2=\frac{1/3 - (p-\rho q)^2}{q^2}\,.\]
 Notice that $\gmin q^2+\gammabar(p-\rho q)^2\geq
 \frac{\gmin}{1+1.5\rho^2}$ if and only if $\gmin/\gammabar\leq f_1$.
 On the other hand, $\gmin q^2+\gammabar(p-\rho q)^2\geq \gammabar/3$ if
 and only if $\gmin/\gammabar\geq f_2$.  Thus, it suffices to show
 $f_1\geq f_2$ for all parameter choices.
First, notice that since $\frac{1}{1+1.5\rho^2}-q^2\geq 0$, we have  $f_1\geq f_2$ whenever
 \[1\geq \left(\frac{1}{3(p-\rho
   q)^2}-1\right)\left(\frac{1}{(1+1.5\rho^2)q^2}-1\right)\,.\]
This is satisfied whenever $3(p-\rho q)^2\geq 1-(1+1.5\rho^2)q^2$.
Expanding and bringing all to the left hand side shows this is true because $(2p-3\rho q)^2\geq 0$.
\end{proof}
%

While we do not currently have a comparison between our
\expref{Theorem}{thm:main} and Corollary~2 of~\cite{jstv04}, which requires
a pointwise bound on $\pi_i^j$, we believe that Corollary~2 of~\cite{jstv04}
is insufficient for our application to permutations.  We expound
upon this in \expref{Remark}{rem:eta}.


\section{Application to permutations}\label{sec:perm}
In this section, we apply the complementary decomposition theorem, 
\expref{Theorem}{thm:productspaces}, to the problem of sampling biased permutations.  
We give the proofs of \expref{Theorem}{pprocess} and \expref{Theorem}{thm:mnGap} which bound the spectral gap of the nearest-neighbor Markov chain $\mn$ and the $k$-particle process Markov chain $\mk$, respectively.  
The section is laid out as follows.  In \expref{Section}{sec:MKprelim}, 
we describe the state spaces and stationary distributions of $\mn$ and $\mk,$ and formally give the definitions of the weakly monotonic 
property for each chain.  Next, in \expref{Section}{sec:mk} we give the proof of 
\expref{Theorem}{pprocess}, and then in \expref{Section}{sec:ortho} we complete the 
proof by proving that our decomposition is 
%
$\epsilon$-orthogonal, where $\epsilon\leq 1/n^2$ %
(\expref{Lemma}{lem:perms}). Finally, in \expref{Section}{sec:mn} we use comparison techniques~\cite{dsc, RT98} to bound the spectral gap of $\mn$ (\expref{Theorem}{thm:mnGap}) using the bound on the spectral gap of $\mk$ (\expref{Theorem}{pprocess}).   

\subsection{Biased permutations and $k$-particle processes}\label{sec:MKprelim}

We begin by giving the formal definition for the nearest neighbor 
Markov chain $\mn$ over biased permutations. Here we are interested in 
the problem of sampling from the symmetric group $S_n$, the 
permutations of $[n]$, where we interpret the elements of $S_n$ as strings
of length $n$ with each symbol in $[n]$ appearing exactly once.  
We are given as input a set of probabilities 
$\cP=\{p_{i,j}\}$ for all $1\leq i,j\leq n$ with $p_{i,j}=1-p_{j,i}$.  
Let $\sigma(i)$ denote the element in position $i$ of $\sigma \in S_n$.

\vspace{.1in}
\noindent  {\bf The Nearest Neighbor Markov chain $\mn$ } 

\vspace{.05in}
\noindent {\tt Starting at any permutation $\sigma^0$, iterate the following:} 
\begin{itemize}
\item At time $t\geq0$, choose a position $1< i\leq n$ uniformly at random in permutation $\sigma^t$.  
\item With probability $p_{\sigma^t(i), \sigma^t(i-1)}/2$, exchange the elements $\sigma^t(i)$ and $\sigma^t(i-1)$ to obtain $\sigma^{t+1}$.
\item Otherwise, do nothing so that $\sigma^{t+1} = \sigma^{t}.$
\end{itemize}

\noindent The chain $\mn$ connects the state space $S_n$ and has the 
following stationary distribution (see, \eg,~\cite{bmrs}): 
$$\pin(\sigma) = \prod_{i < j: \sigma(i) > 
\sigma(j)}\frac{p_{\sigma(i),\sigma(j)}}{p_{\sigma(j),\sigma(i)}} 
\Zn^{-1} = \prod_{i < j: \sigma(i) > \sigma(j)}q_{\sigma(i), \sigma(j)} 
\Zn^{-1} $$ where $\Zn$ is a normalizing constant  and $q_{\sigma(i), 
\sigma(j)}=\frac{p_{\sigma(i),\sigma(j)}}{p_{\sigma(j),\sigma(i)}}.$  

Our central question is under what conditions does $\mn$ mix in 
polynomial time. %
Fill~\cite{F03}    %
introduced monotonicity 
conditions under which he conjectured $\mn$ would be rapidly mixing. In 
our analysis of $\mn$ we use the weakly monotonic condition that appears
in~\cite{bmrs,MS}: 

\begin{definition}[\cite{bmrs}]\label{mono}The set $\p$ is weakly monotonic if properties 1 and either 2 or 3 are satisfied.
\begin{enumerate}
\item $p_{i,j} \geq 1/2$ for all $1 \leq i < j \leq n,$ and
\item $p_{i,j+1} \geq p_{i,j}$ for all $1 \leq i < j \leq n-1$ or
\item $p_{i-1,j} \geq p_{i,j}$ for all $2 \leq i < j \leq n.$
\end{enumerate}
\end{definition}

We consider the special case of $k$-classes where $[n]$ is partitioned 
into $k$ classes $\mathcal{C}_1,\mathcal{C}_2,\ldots, \mathcal{C}_k$, 
and assume elements in class $\mathcal{C}_i$ interact with elements in 
class $\mathcal{C}_j$ with the same probability.  That is, if 
$i_1,i_2\in \mathcal{C}_i$ and $j_1,j_2\in \mathcal{C}_j$ then 
$p_{i_1,j_1} = p_{i_2,j_2}$.  In this case we define $\pbar_{i,j}$ to 
be this shared probability for classes $\mathcal{C}_i$ and 
$\mathcal{C}_j$ (the bar indicates that we have reindexed the set of 
probabilities by the classes) and we say that $\p$ forms a 
\emph{k-class}.  Note that $\pbar_{i,i}$ is assumed to be $1/2$, so 
that $\mn$ swaps elements within the same class with probability $1/2$.  
Define $\cPbar=\{\pbar_{i,j}\}$ as the set of probabilities over pairs 
of classes $\mathcal{C}_i$ and $\mathcal{C}_j$ where $i,j\leq k$.  In 
this case, we say $\cP$ forms a \emph{$k$-class}.  When $k=n$, the 
$k$-class assumption does not lose any generality, but this structure 
allows us to simplify the problem by considering $k<n$, as was done 
in~\cite{MS,HW16}. 

Define $\qbar_{i,j} =  \pbar_{i,j}/\pbar_{j,i}$ to be the \emph{bias} 
towards having a particle of type $i$ ahead of a particle of type $j$.  
We say that $\cP$ is \emph{bounded} if there exists a constant $q>1$ 
such that $\qbar_{i,j}\geq q$ for all $i<j\leq k$.  The constant $q$ is 
called the \emph{minimum bias}. 

The chain $\mn$ samples over $S_n$ using these probabilities, and in 
particular the order of elements within each class approaches the 
uniform distribution.  The spectral gap of this uniform sampling is 
well-understood and may be analyzed separately (see 
\expref{Section}{sec:mn}).  In order to isolate the biased moves, we define 
a new Markov chain $\mk$ that eliminates swaps within each class.  As 
$\mk$ maintains a fixed order on particles within each class, it makes 
sense to relabel each element of $[n]$ by the index of the class it is 
in.  That is, we let $c_i = |\mathcal{C}_i|$ and we consider a linear 
array of length $n$ with $c_i$ particles labeled $i$ for each $1\leq 
i\leq k$.  We call this a $k$-particle system for the given 
set~$\{c_i\}$, and the Markov chain $\mk$ is called a $k$-particle 
process.  We view the new state space as the set $\Omega$ of 
$k$-particle systems for~$\{c_i\}$.


The Markov chain $\mk$ over $k$-particle systems also allows certain 
non-adjacent transpositions.  In particular, we let a particle of type 
$i$ and a particle of type $j$ swap across particles of type less than 
$i$ and $j$.  More formally, the chain $\mk$ is defined as follows.


\vspace{.1in}
\noindent  {\bf The particle process Markov chain $\mk$} 

\vspace{.05in}
\noindent {\tt Starting at any $k$-particle system $\sigma^0$, iterate the following:} 
\begin{itemize}
\item At time $t,$ choose a position $1\leq i \leq n$ and direction 
$d \in \{L, R\}$ uniformly at random.  
\item If $d = L$, find the largest $j$ less than $i$ with $\sigma^t(j) \geq \sigma^t(i)$ (if one exists).  If $\sigma^t(j) > \sigma^t(i),$ then 
with probability $1/2$,
exchange $\sigma^t(i)$ and $\sigma^t(j)$ to obtain $\sigma^{t+1}.$  
\item If $d = R$, find the smallest $j$ with  $j>i$ and $\sigma^t(j) \geq \sigma^t(i)$ (if one exists).  If $\sigma^t(j) > \sigma^t(i),$ then exchange $\sigma^t(i)$ and $\sigma^t(j)$ to obtain $\sigma^{t+1}$ with probability 
$$\frac{1}{2}~ \qbar_{\sigma_t(j),\sigma_t(i)}  \prod_{i<l<j} \qbar_{\sigma_t(j),\sigma_t(l)}\qbar_{\sigma_t(l),\sigma_t(i)}\,\,.$$


\item With all remaining probability, $\sigma^{t+1} = \sigma^t$.
\end{itemize}

\noindent The chain $\mk$  connects the space $\Omega$ and has the stationary distribution (see, \eg,~\cite{bmrs})
$$\pi(\sigma) = \prod_{i < j: \sigma(i) > \sigma(j)}\frac{\pbar_{\sigma(i),\sigma(j)}}{\pbar_{\sigma(j),\sigma(i)}} Z^{-1} = \prod_{i < j: \sigma(i) > \sigma(j)}\qbar_{\sigma(i), \sigma(j)} Z^{-1} $$ where $Z$ is a normalizing constant and $\qbar_{\sigma(i), \sigma(j)}=\frac{\pbar_{\sigma(i),\sigma(j)}}{\pbar_{\sigma(j),\sigma(i)}}.$  

We adapt the definition of weakly monotonic to the setting of $k$-particle systems as follows:
 
\begin{definition}[\cite{bmrs}]\label{mono}The set $\cPbar$ is weakly monotonic if properties 1 and either 2 or 3 are satisfied.
\begin{enumerate}
\item $\pbar_{i,j} \geq 1/2$ for all $1 \leq i < j \leq k,$ and
\item $\pbar_{i,j+1} \geq \pbar_{i,j}$ for all $1 \leq i < j \leq k-1$ or
\item $\pbar_{i-1,j} \geq \pbar_{i,j}$ for all $2 \leq i < j \leq k.$
\end{enumerate}
\end{definition}
\noindent As in~\cite{MS}, we will assume that property (2) holds.  If
instead property (3) holds, then as described in~\cite{MS} we would
modify $\mk$ (and $\mtk$ defined below) to allow swaps between
elements of different particle types across elements whose particle
types are larger (instead of smaller) and modify the induction so
that at each step $\sigma_{i}$ restricts the location of particles
larger than $i$ (instead of smaller).


%
%
\subsection{Details of the decomposition and proof of \expref{Theorem}{pprocess}}
\label{sec:mk}

In this section, we present the proof of \expref{Theorem}{pprocess}, which bounds 
the mixing time of the particle process chain $\mk.$  

\begin{thm:mk}
If the probabilities $\cP$ are weakly monotonic and form a bounded 
$k$-class with  $c_i\geq 2\psize$ for all $1\leq i \leq k$, 
then the spectral gap $\gamma$ of the chain $\mk$ satisfies
\abovedisplayskip=3pt
\belowdisplayskip=3pt
$\gamma =\Omega(n^{-2}).$
\end{thm:mk}


Our proof uses the same inductive technique
as~\cite{MS}, where at each level of the induction we fix the
locations of particles in one less particle class.  
Recalling notation from \expref{Section}{sec:intro}, 
let $\mathcal{C}_{<= i} = \mathcal{C} \cup \mathcal{C}_{< i}$ 
be the set of particles with type less than or 
equal to $i$ (\ie,  particles of type $1, 2, \ldots, i$), and similarly,
$\mathcal{C}_{>= i} = \mathcal{C} \cup \mathcal{C}_{> i}$.  For $i\geq 0$, 
let $\sigma_i$ represent a fixed location of 
the particles in $\mathcal{C}_{<= i}$ ($\sigma_0$ represents 
no restriction); for example, in \expref{Figure}{fig:product}, we set 
$\sigma_{2} = 12\_1\_2\_1\_\_$, where ``$\_$'' represents  
locations that can be filled with particles in $\mathcal{C}_{>= 3}$. 
We will consider the chain $\m_{\sigma_{i}}$ whose state space 
$\Omega_{\sigma_i}$ is the set of all $k$-particle systems  $\sigma$ 
where the location of particles in $\mathcal{C}_{<= i}$ 
 are consistent with $\sigma_i$. 
The moves of $\m_{\sigma_{i}}$ are those moves from $\mk$ that 
do not involve particles in $\mathcal{C}_{<= i}$.
We prove by induction that 
$\m_{\sigma_{i}}$ has spectral gap $\Omega(n^{-2}(1-1/n)^{2(k-2-i)})$ 
for all choices of $\sigma_i$. To be clear, we assume that the spectral gap of
$\m_{\sigma_{i}}$ are bounded for all $\sigma_i$ by induction, and then 
prove our bound on the spectral gap of $\m_{\sigma_{i-1}}$.


 \begin{figure}
   \centering
   \begin{minipage}[c]{.5\textwidth}
     \centering
   \includegraphics[scale=.18]{product}
   \put(-30,98){$\Omegahat_{a}=\Omega_{\sigma_i}$}
   \put(-28,88){$\downarrow$}
   \put(-28,58){$\sigma $}
  \put(0,58){$\leftarrow\Omegatilde_b$}
   \end{minipage}\hfill
   \begin{minipage}[c]{.5\textwidth}
     \centering
     \begin{tabular}{rl}
     $\sigma_2$ &$=~~ 12\_1\_2\_1\_\ \_$\\
     $\sigma_3$ &$=~~
     12{\color{red}3}1{\color{red}\_}2{\color{red}\_}1{\color{red}3\_}$\\
     $a$ &$= ~~~ {\color{red} 3\_\ \_3\_}$\\
     $b$ &$= ~~~ {\color{blue} 454}$\\
     $\sigma$ &$=~~
     12{\color{red}3}1{\color{blue}4}2{\color{blue}5}1{\color{red}3}{\color{blue}4}$     
     \end{tabular}
     \end{minipage}
   \caption{The state space $\Omega_{\sigma_{i-1}}$ decomposed, with $i=3$.}
   \label{fig:product}
 \end{figure}

To start, we show that $\Omega_{\sigma_{i-1}}$ is a product space,
which is required to apply \expref{Theorem}{thm:productspaces}.
Let $A$ consist of all 2-particle systems with $c_i = |\mathcal{C}_i|$ particles of type $i$ 
and $\sum_{j=i+1}^k c_j$ particles of type ``$\_$''.  Let 
$B$ consist of all $k-i$ particle systems with $c_j = |\mathcal{C}_j|$ particles of type $j$
for all $i+1 \leq j \leq k$. For example, using our running example where $\sigma_{2} = 12\_1\_2\_1\_\_,$ an example $a \in A$ is $\_3\_3\_$ and $b \in B$ is $544.$ See \expref{Figure}{fig:product} for an additional example.  Our goal is to show that the set of 
permutations $\sigma$ consistent with $\sigma_{i-1}$ on particles in 
$\mathcal{C}_{< i}$ is in bijection with $A \times B$.  To this end, we 
can write $\sigma=(a,b)$, where $a \in A$ is the 2-particle system 
obtained from $\sigma_i$ by removing particles in $\mathcal{C}_{< i}$ 
(see \expref{Figure}{fig:product}), as those particles are in a fixed 
position for all of $\Omega_{\sigma_{i-1}}$.
 Next, define $b\in B$ to be the restriction of $\sigma$ to 
particles in $\mathcal{C}_{>i}$.  For the other direction, given any 
$(a,b)$ pair, it is clear that there is a unique 
$\sigma\in\Omega_{\sigma_{i-1}}$ corresponding to that pair.  Returning to our example, given $\sigma_{2} = 12\_1\_2\_1\_\_,$ $a=\_3\_3\_,$ and $b =544,$ then the unique $\sigma \in \Omega_{\sigma_{2}}$ is $1251324134.$

\begin{figure}[h]
  \centering
   \begin{tikzpicture}[scale=.6]
     \draw [fill=lightgray,thin,lightgray] (0,0) rectangle (1,3);
     \draw [fill=lightgray,thin,lightgray] (1,0) rectangle (2,2);
     \draw [fill=lightgray,thin,lightgray] (2,0) rectangle (4,1);
              \draw [fill=gray,thin,gray] (2,1) rectangle (3,2);
     \draw [help lines] (0,0) grid (4,3);
     \draw [ultra thick] (0,3) -- (1,3) -- (1,2) -- (3,2) -- (3,1) -- (4,1) -- (4,0);
     \node[draw] at (1.95,-.55) {$i~\_~i~i~\_~i~\_$};
   \end{tikzpicture}  
\hspace{2em}
    \begin{tikzpicture}[scale=.6]
     \draw [fill=lightgray,thin,lightgray] (0,0) rectangle (1,3);
     \draw [fill=lightgray,thin,lightgray] (1,0) rectangle (2,2);
     \draw [fill=lightgray,thin,lightgray] (2,0) rectangle (4,1);
         \draw [fill=gray,thin,gray] (2,1) rectangle (3,2);
     \draw [help lines] (0,0) grid (4,3);
     \draw [ultra thick] (0,3) -- (1,3) -- (1,2) -- (2,2) -- (2,1) -- (4,1) -- (4,0);
          \node[draw] at (1.95,-.55) {$i~\_~i~\_~i~i~\_$};
   \end{tikzpicture}
\caption{An exclusion process on staircase walks operates by adding or removing a square. }
\label{fig:staircase}
\end{figure}


We next describe the decomposition.  Note that the moves of 
$\m_{\sigma_i}$ fix an $a \in A$ and 
perform $(j_1,j_2)$ transpositions, where $j_1,j_2>i$; \ie,  they 
operate exclusively on $B$. Thus, the Markov chain $\m_{\sigma_i}$ is a 
restriction of $\m_{\sigma_{i-1}}$ with state space 
$\Omegahat_{a}$.
On the other hand, the remaining moves of $\m_{\sigma_{i-1}}$ 
are $(i,j)$ transpositions for $j>i$.  These are the complementary 
restrictions; these moves fix a $b \in B$ and operate on $A$, so we 
label the state space of this Markov chain $\Omegatilde_b$.  As these 
moves fix the relative order of all particles in $\mathcal{C}_{>i}$, 
the complementary restriction chains can be seen as a 2-particle process on particles in $\mathcal{C}_{i}$  with particles in $\mathcal{C}_{> i}$.  This process is easily seen as in bijection with staircase walks by mapping each particle in $\mathcal{C}_{i}$ to a step right and each particle in $\mathcal{C}_{> i}$ to a step down, as in \expref{Figure}{fig:staircase}.  
In~\cite{MS}, they define a \emph{bounded generalized biased exclusion process}  as a processes that operate on 
staircase walks as in \expref{Figure}{fig:staircase}, where every square 
has a different bias but they are all bounded by some $q$.  More formally, an exclusion process is a process on $n_1$ 1's and $n_0$ 0's which occupy linear positions: $1, \ldots, n_0 + n_1.$. The set of all distinct linear orderings of these particles are called \emph{2-particle systems}.  In the generalized setting, for a configuration $\sigma$, the probability $p_{\sigma, i}$ of swapping two particles at position $i$ and $i+1$  can depend on both the current ordering and the particles being exchanged.  The process is \emph{bounded} if there exists  a constant $\gamma > 1$ such that for all configurations $\sigma$ if $\sigma(i)  =1$ and $\sigma(i+1) = 0$ and $\tau$ is obtained from $\sigma$ by swapping particles $\sigma(i)$ and $\sigma(i+1),$ then $p_{\sigma,i}/p_{\tau_i}\geq q.$  Bounded generalized biased exclusion 
processes were analyzed in~\cite{MS}. 

\begin{theorem}[\cite{MS}]\label{boundedBias} Let $\me$ be a bounded 
generalized exclusion process on $n_1$ 1's and $n_0$ 0's.  The spectral 
gap of $\me$ is $\Omega((n_0n_1)^{-1})$.
\end{theorem}

Next we prove the following lemma.  

\begin{lemma}\label{lem:dual}
The complementary restrictions at each level of the induction are 
bounded generalized biased exclusion processes with spectral gap 
$\Omega(n^{-2})$.
\end{lemma}
\begin{proof}
Since $b$ is fixed, each particle in $b$ corresponds to a given row of 
squares in \expref{Figure}{fig:staircase}.  Each column in the figure 
corresponds to a particular element of type $i$ in $a$.  The bias of 
swapping an element $i$ with a particular ``$\_$'' depends only on the 
position of those two elements in the array $\sigma$, and not on the 
relative order of all other elements in $\mathcal{C}_{> i-1}$.  In 
this way, that bias corresponds to a particular square in 
\expref{Figure}{fig:staircase}.  Elements in $\mathcal{C}_{<i}$ change the 
bias of all squares along a given diagonal in 
\expref{Figure}{fig:staircase}.  For example, let $\sigma = 
3243\underline{3}1\underline{5}324$ and suppose $i=3$.  Then 
$a=3\_33\_3\_$ and $b=454$.  The highlighted square in 
\expref{Figure}{fig:staircase} corresponds to swapping the 3 and 5 in 
positions 5 and 7, respectively.  The probability of making that move 
is $\qbar_{5,3}\qbar_{5,1}/2\qbar_{3,1}$.  Thus, each complementary 
restriction chain can be viewed as a generalized exclusion process 
acting on $a\in A$.  The minimum bias is $q:=\min_{i<j} \qbar_{i,j}$, 
which we have assumed is a constant bigger than 1. We use the following 
result from~\cite{MS}.


\noindent Since the probabilities $\cPbar$ are weakly monotonic 
(specifically, they satisfy condition (1) of weak monotonicity) and 
bounded, the exclusion process involving the particles in 
$\mathcal{C}_{i}$ and the particles in $\mathcal{C}_{<i}$ is a 
bounded generalized exclusion process and we can apply 
\expref{Theorem}{boundedBias}.  There are $c_i$ particles of type $i$, 
$\sum_{j=i+1}^k c_j < n$ particles of type ``$\_$'', and the moves are 
selected with probability $c_i/4n$.  Applying 
\expref{Theorem}{boundedBias} shows that the spectral gap of each 
complementary restriction chain is 
$\Omega(\frac{c_{i}}{n}\frac{1}{nc_{i}})=\Omega(n^{-2})$. \end{proof}


The chain $\m_{\sigma_{i-1}}$ is not the direct product of the chains 
on $A$ and $B$ because, \eg, for $(a,b) \in A\times B$, 
$P((a,b),(a',b))$ depends on $b$. However, in \expref{Section}{sec:ortho} 
we prove \expref{Lemma}{lem:perms} which states that the above 
decomposition is 
%
$\epsilon$-orthogonal where $\epsilon\leq 1/n^2$.%
Specifically, \expref{Lemma}{lem:perms} says that if the probabilities $\p$ are weakly monotonic and bounded with 
$c_i\geq 2\psize$ for all $1\leq i \leq k$, then at each step of 
the induction, %
$\epsilon\leq 1/n^2$. %
Here, we use 
\expref{Lemma}{lem:perms} to complete the proof of \expref{Theorem}{pprocess} 
which says that if the probabilities $\cPbar$ are weakly monotonic and 
bounded, and  $c_i\geq 2\psize$ for $1\leq i \leq k$ then the spectral 
gap $\gamma$ 
of the chain $\mk$ satisfies $\gamma =\Omega(n^{-2})$.  Note that
$\cP$ being weakly monotonic implies that $\cPbar$ is as well.
%

%
\begin{proof}[Proof of \expref{Theorem}{pprocess}]
At each step of the induction, we apply the complementary decomposition 
theorem (\expref{Theorem}{thm:productspaces}). The restrictions of each 
decomposition will be rapidly mixing by induction and the complementary 
restriction chains are bounded generalized exclusion processes, by 
\expref{Lemma}{lem:dual}.  The base case is $i=k-2$ and the final 
decomposition is $i=0$.
 
We begin with our base case, $i=k-2$.  
Let $\sigma_{k-2}$ be any fixed location of the particles in $\mathcal{C}_{<= k-2}.$ 
The Markov chain $\m_{\sigma_{k-2}}$ rejects all moves of $\mk$ unless 
they exchange a particle in $\mathcal{C}_{k-1}$ with a particle in 
$\mathcal{C}_{k}$,
making it a generalized exclusion process. Thus, 
$\m_{\sigma_{k-2}}$ is a bounded generalized exclusion process slowed 
down by a factor of $c_{k-1}/4n$, as this is the probability that 
these transitions are chosen.  By \expref{Theorem}{boundedBias} for any 
such $\sigma_{k-2}$, $\m_{\sigma_{k-2}}$ has spectral gap 
$\Omega(\frac{c_{k-1}}{n}\frac{1}{n c_{k-1}})=\Omega(n^{-2})$. 

We assume by induction the mixing time bound holds for all 
$\m_{\sigma_i}$ for some $i\leq k-2$, and we will use this result to 
prove that our mixing time bound holds for all $\m_{\sigma_{i-1}}$. Let 
$\sigma_{i-1}$ represent any fixed choice of locations for all 
particles in $\mathcal{C}_{< i}$.  In order to bound the spectral gap 
$\gamma( \m_{\sigma_{i-1}})$ of the chain $\m_{\sigma_{i-1}}$ we will 
apply \expref{Theorem}{thm:productspaces}.  Given any $\sigma_i$ that is 
consistent with $\sigma_{i-1}$ (\ie,  they agree on the locations of all 
particles in $\mathcal{C}_{<= i}$), the Markov chain $\m_{\sigma_i}$ 
will be a restriction Markov chain of $\m_{\sigma_{i-1}}$.  By 
induction, we have that the spectral gap $\gamma(\m_{\sigma_i})= 
\Omega\left(n^{-2}\left(1-\frac{1}{n}\right)^{2(k-2-i)}\right).$ 
\expref{Lemma}{lem:dual} shows that the complementary restrictions have 
minimum gap $\Omega(n^{-2}).$ By \expref{Lemma}{lem:perms}, we know that 
the decomposition is %
$\epsilon$-orthogonal, with $\epsilon\leq 1/n^2$. %
Combining these with 
\expref{Theorem}{thm:productspaces} implies $\gamma(\m_{\sigma_{i-1}})= 
\Omega\left(n^{-2}\left(1-\frac{1}{n}\right)^{2(k-2-i)+2}\right)$. 
Substituting $i = 0$ gives the desired theorem
%
\[
\gamma(\mk) = 
\Omega\left(n^{-2}\left(1-\frac{1}{n}\right)^{2k-2}\right) = 
\Omega\left(n^{-2}\left(1-\frac{1}{n}\right)^{n}\right) = 
\Omega\left(n^{-2}\right)\,.\qedhere
%
\]
\end{proof}
%

\begin{remark}\label{rem:eta}

It is worth pointing out that this decomposition of $\m_{\sigma_{i-1}}$ 
does not satisfy the 
regularity conditions of~\cite{jstv04} needed to obtain a better bound.  
For any $\upsilon\in \Omegahat_j$, define 
\[\pihat_j^{j'}(\upsilon)=\pi_j(\upsilon)\frac{\sum_{\upsilon'\in\Omegahat_{j'}}P(\upsilon,\upsilon')}{\Pbar(j,j')}\,.\] 
We need to bound $\pihat_j^{j'}(\upsilon)/\pi_j(\upsilon)$ for any 
$j,j',$ and $\upsilon\in\Omegahat_j$.   For example, let $\sigma_{2} = 
12\_11111\_2\_1\_~\_$ and $\sigma_3 = 12311111\_231\_~\_$. Notice that 
the two permutations $\upsilon_1=12311111423156$ and 
$\upsilon_2=12311111523146$ are in the same restriction $\Omegahat_j$ 
(\ie,  they are both consistent with $\sigma_3$).  They each have a 
single move to $\Omegahat_{j'}$: the move of swapping the first 3 with 
the 4 (in the case of $\upsilon_1$) or 5 (in the case of $\upsilon_2$).  
However, the probability of these moves differ by a factor of 
$(\qbar_{4,3}/\qbar_{5,3})(\qbar_{4,1}/\qbar_{5,1})^5,$ as there are five 1's between 
them.  In principle, there could be order $n$ smaller numbers between 
the two numbers we are swapping.  Thus, 
$\pihat_j^{j'}(\upsilon)/\pi_j(\upsilon)$ cannot be uniformly bounded 
to within $1\pm \eta$ unless $\eta$ is exponentially large. 

\end{remark}

%
%
\subsection{Bounding the orthogonality of the decomposition.}\label{sec:ortho}  %

In order to complete the proof of \expref{Theorem}{pprocess}, it remains to prove
\expref{Lemma}{lem:perms}, which states that %
%
at each step of the induction,
the value of $\epsilon$ is bounded by $1/n^2$. %

\begin{lemma}\label{lem:perms}
If the probabilities $\p$ are weakly monotonic and bounded with 
$c_i\geq 2\psize$ for all $1\leq i \leq k$, then at each step of 
the induction %
the quantity $\epsilon$ defined in \expeqref{equation}{eq:eps} 
satisfies $\epsilon\leq 1/n^2$. %
\end{lemma}

Before proceeding with the proof, 
recall that the chain $\m_{\sigma_{i-1}}$ is not the direct product of 
the chains on $A$ and $B$ because, \eg, for $(a,b) \in A\times B$, 
$P((a,b),(a',b))$ depends on $b$. However, we show that the 
decomposition given in \expref{Section}{sec:mk} is $1/n^2$-orthogonal by 
bounding $r(a,b) = \pi(a,b)/(\pi(a)\pi(b))$.  Recall that at each step 
of the induction, $a$ is a 2-particle system obtained from $\sigma_i$ 
by removing all particles in $\mathcal{C}_{< i}$, $c_i = 
|\mathcal{C}_i|$, and there are $\sum^{k}_{j=i+1}c_j$ particles of type 
``$\_$''.  Also recall that $b$ is the $(k-i)$-particle system 
consisting of the particles in $\mathcal{C}_{>i}$, and the location of 
the particles in $\mathcal{C}_{< i-1}$ are fixed throughout this step 
of the induction (see \expref{Figure}{fig:product}).  

We define $a$ to be ``good'' when it has fewer than $N^*=C_q\log n$ 
inversions; otherwise $a$ is ``bad.''  Similarly, we define $b$ to be 
``good'' when it has fewer than $N^*$ inversions involving particles in 
$\mathcal{C}_{i+1}$; otherwise it is ``bad.''  Thus, as 
$|\mathcal{C}_{i+1}| \geq 2N^*$, $(a,b)$ has no inversions between $i$ 
and $j$ for $j>i+1$ when $a$ and $b$ are both good.  For such pairs, 
$r(a,b)$ is very close to 1.  We will show that almost all weight 
contributing to $\pi$ comes from pairs $(a,b)$ where $a$ and $b$ 
are good.  For all other pairs we show 
$r(a,b)\pi(\Omegahat_a\cap\Omegatilde_{b})$ is small.  

By viewing $b$ 
as a staircase walk on $\mathcal{C}_{i+1}$ and $\mathcal{C}_{> i+1}$, 
we see that for either $a$ or $b$, the 
probability it is bad is smaller than the weighted sum of all biased 
exclusion processes with more than $N^*$ inversions (equivalently, area 
$N^*$ under the curve).  We start with the following lemma.

\begin{lemma}\label{lem:partitions}
For any biased exclusion process with minimum bias constant $q > 1$, 
the total weight of staircase walks with area larger than $\psize$ satisfies 
\[\sum_{\sigma:\area(\sigma)\geq \psize}q^{-\area(\sigma)} \leq \frac{1}{6n^2}\,.\]
\end{lemma}

\begin{proof}


The number of staircase walks with area $\area$ under the curve is 
precisely the partition number $p(\area)$.  By a theorem of 
Erd\H{o}s~\cite{Erdos42}, $p(\area) < \eee^{\pi\sqrt{2\area/3}}$,
and as $\eee^{\pi\sqrt{2/3}}<14$, we have $p(\area) < 14^{\sqrt{\area}}$.
We can use this bound on $p(\area)$ to bound 
the weight of a set of staircase walks, where each inversion 
contributes a factor of $1/q$ to the weight of the staircase walk 
(recall $q:=\min_{i<j} \qbar_{i,j}$).  
Define 
\[\psize = \max\left\{\frac{\log (6n^2)+\log((1+q)/(q-1))}{\log((1+q)/2)}, \frac{\log^2(14)}{\log^2(2q/(1+q))}\right\} = \Theta(\log n)\,.\]
  Let $N(\sigma)$ be the number of inversions in a particular staircase walk $\sigma$.  
  For all $\area\geq
  \frac{\log^2(14)}{\log(2q/(1+q))}$, we have $14^{\sqrt{\area}}q^{-\area}\leq (2/(1+q))^\area.$
  Therefore 
\[  \sum_{\sigma:\area(\sigma)\geq \psize}q^{-\area(\sigma)}\leq \sum_{\area\geq
    \psize}p(\area)q^{-\area} \leq \sum_{\area\geq \psize}14^{\sqrt{\area}}q^{-\area}
%
  \leq \sum_{\area\geq \psize}\left(\frac{2}{1+q}\right)^{\area}  %
     =\left(\frac{2}{1+q}\right)^{\psize}\frac{1}{1-\frac{2}{1+q}}
       \leq \frac{1}{6n^2}\]
by the choice of $\psize$.        
\end{proof}

%
\subsubsection{High-level summary}  %
 We are now ready to use \expref{Lemma}{lem:partitions} to prove 
that  %
$\epsilon\leq 1/n^2$. %
For improved 
readability, we illustrate the main ideas of the proof here 
in the simplified $k=3$ case before proving \expref{Lemma}{lem:perms}. 
In this case, there is no recursion but 
instead just a single application of the decomposition theorem.  The 
restriction chains of $\m=\m_{\sigma_0}$ are the set 
$\{\m_{\sigma_1}\}$, which 
fix all elements in class $\mathcal{C}_1$.  
The stationary distribution of $\m$ is
\begin{equation}\label{eq:statdist}
\pi(\sigma) = Z^{-1} \prod_{i < j:\atop \sigma(i) > \sigma(j)}\qbar_{\sigma(i), \sigma(j)},
\end{equation}
where $Z$ is a normalizing constant.

Let $w(a)$ and $ w(b)$ be the parts of this product that depend only on 
$a$ and only on $b$, respectively, and let $w(a,b)$ be a correction 
factor that depends on both $a$ and $b$.  Let $t_{1,3}$ (respectively, 
$t_{1,2}$ and $t_{2,3}$) denote the number of inversions in $\sigma$ 
between a 1 and a 3 (respectively a 1 and a 2 and a 2 and a 3). For 
example, let $\sigma = 111221312323$, which has stationary probability 
$Z^{-1}(\qbar_{1,2})^4(\qbar_{2,3})^3(\qbar_{1,3})$.  Then $a=111\_\ \_1\_1\_\ \_\ \_\ \_$ 
$b=2232323$.  From $b$ we find that $t_{2,3}=3$ and 
$w(b)=(\qbar_{2,3})^3$; more generally, define $w(b)$ to be the 
product~$(\qbar_{2,3})^{t_{2,3}}$. From $a$, we can see that there are five 
inversions involving 1, but the number of those that are inversions 
with a $3$ versus a $2$ depends on $b$ as well. Ignoring this for a 
moment, we define $w(a)$ to be the product $(\qbar_{1,2})^{t_{1,2}+t_{1,3}}$. 
In our example, $w(a)=(\qbar_{1,2})^5$.  Since we have made the false 
assumption that there were no inversions between a 1 and a 3 in 
$\sigma$, we need a correction factor 
$w(a,b)=(\qbar_{1,3}/\qbar_{1,2})^{t_{1,3}}$.
With these definitions, it is clear that $\pi(\sigma)=Z^{-1}w(a)w(b)w(a,b)$.


A key idea in the proof of \expref{Lemma}{lem:perms} is that if $a$ and $b$ are both good, then 
$t_{1,3}=0$---indeed, the total number of inversions is less than 
$2N^*$ and the number of 2's is at least $2N^*$---and thus the 
correction factor satisfies $w(a,b)=1$, implying $\pi(a,b)\approx \pi(a)\pi(b)$. 
Moreover, the probability that $a$ or $b$ is bad is very small, so 
these pairs $(a,b)$ do not contribute much to the sum in 
\expeqref{equation}{eq:eps}.  

\subsubsection{Proof of \expref{Lemma}{lem:perms}}
We may now prove \expref{Lemma}{lem:perms} in its full generality.\\

%
\begin{proof}[Proof of \expref{Lemma}{lem:perms}]   %
Assume that at this step in the induction, all particles in 
$\mathcal{C}_{< i}$ are fixed in the same position in all $k$-particle 
systems according to some $\sigma_{i-1}.$  The stationary distribution 
of the chain $\m_{\sigma_{i-1}}$ is $$\pi_{\sigma_{i-1}}(\sigma) =  
\prod_{i < j: \sigma(i) > \sigma(j)}\qbar_{\sigma(i), 
\sigma(j)}Z_{\sigma_{i-1}}^{-1}$$  
where $Z_{\sigma_{i-1}}$ is the normalizing constant $\sum_{\sigma \in 
\Omega_{\sigma_{i-1}}} \pi(\sigma),$ the set $\Omega_{\sigma_{i-1}}$ 
contains the $k$-particle systems consistent with $\sigma_{i-1},$ and 
$\qbar_{\sigma(i), 
\sigma(j)}=\frac{\pbar_{\sigma(j),\sigma(i)}}{\pbar_{\sigma(j),\sigma(i)}}.$  
For ease of notation, throughout the remainder of this section we will 
let $\pi = \pi_{\sigma_{i-1}}$ and $Z =Z_{\sigma_{i-1}}$.

Let $a^*$ and $b^*$ be the highest weight elements in $A$ and $B$, 
respectively.  Using these definitions, the $k$-particle system 
$(a^*,b^*)$ has the particles in $\mathcal{C}_{< i}$ fixed according to 
$\sigma_{i-1}$ and all other higher particles in sorted order.   In our 
example, $b^* = 445,$ $(a,b^*) = 1231423145,$ $a^* = 33\_\ \_\ \_,$ 
$(a^*,b) = 1231324154,$ and $(a^*,b^*) = 1231324145.$ 

Next, we will decompose the product over inversions in 
\expeqref{equation}{eq:statdist} into several quantities. Notice that 
$\pi(a^*,b^*)$ is the product over inversions that are in every 
$\sigma\in\Omega$, normalized by $Z_{\sigma_{i-1}}$.  Define $w(a) = 
\pi(a,b^*)/\pi(a^*,b^*).$  This is the product over inversions that are 
in every $\sigma\in\Omegahat_a$ that are not in every 
$\sigma\in\Omega$. Similarly, $w(b) = \pi(a^*,b)/\pi(a^*,b^*)$ is the 
product over inversions in every $\sigma\in\Omegatilde_b$, and $w(a,b) 
= \pi(a,b) \pi(a^*,b^*)/ (\pi(a,b^*)\pi(a^*,b))$ is the product over 
inversions in $\sigma=(a,b)$ beyond those that are required by being in 
$\Omegahat_a$ and $\Omegatilde_b$.  From these definitions, it is clear 
that $\pi(a,b) = w(a)w(b) w(a,b)\pi(a^*,b^*).$  We will prove in the 
following lemma that if both $a$ and $b$ are good then $w(a,b) = 1;$ 
\ie,  that the weight of $(a,b)$ is determined entirely by being in 
$\Omegahat_a$ and $\Omegatilde_b$.

\begin{lemma}\label{lem:good}
If $a$ and $b$ are both good then $w(a,b) = 1.$
\end{lemma}
\begin{proof}  
From the definitions, it is sufficient to show that $\pi(a,b) = 
\pi(a,b^*)\pi(a^*,b)/\pi(a^*,b^*).$ First consider the inversions 
$(j,l)$ for $j,l<i$.  These are present in every term---$\pi(a,b)$, 
$\pi(a,b^*)$, $\pi(a^*,b)$, and $\pi(a^*,b^*)$---and thus cancel. The 
inversions $(i+1,i)$ are exactly the same in $\pi(a,b)$ and 
$\pi(a,b^*)$, and there are none in $\pi(a^*,b),$ and $\pi(a^*,b^*);$ 
thus these also cancel.  Similarly, inversions $(j,l)$ for  $j\geq i+1,l 
\geq i$ are exactly the same in $\pi(a,b)$ and $\pi(a^*,b),$ and there 
are none in $\pi(a,b^*),$ and $\pi(a^*,b^*)$; thus these cancel. 
Next consider the inversions $(j,l)$ for  $j>i+1,l < i$.  Since there 
are no $(j,i)$ inversions, these are the same in $\pi(a,b)$ and 
$\pi(a^*,b)$ and the same in $\pi(a^*,b^*)$ and $\pi(a, b^*).$ 
Similarly, the $(i,j)$ for $j < i$ inversions are the same in 
$\pi(a^*,b)$ and $\pi(a^*,b^*)$ and the same in $\pi(a,b)$ and 
$\pi(a,b^*)$.  

Finally, the $(i+1,j)$ for $j<i$ inversions are the most complicated. 
Assume the particles in $\mathcal{C}_{i+1}$ are numbered and in order 
in all 4 $k$-particle systems.  Consider any particular particle in 
$\mathcal{C}_{i+1}$ and look at its position in $\pi(a,b)$.  If it is 
to the left of any particle in $\mathcal{C}_{i}$ (\ie,  it is involved 
in an $(i+1,i)$ inversion) then it must come before all particles in 
$\mathcal{C}_{> i+1}$ and it will be in the same position in $\pi(a,b)$ 
and $\pi(a,b^*)$ and the same positions (which may be different 
positions) in $\pi(a^*,b)$ and $\pi(a^*,b^*)$; thus any $(i+1,j)$ 
inversions will cancel.  If not, then it comes after all particles in 
$\mathcal{C}_{i}$ and it will be in the same position in $\pi(a,b)$ and 
$\pi(a^*,b)$ and again the same positions in $\pi(a,b^*)$ and 
$\pi(a^*,b^*)$.  

As there are no $(j,i)$ inversions for $j > i+1$, we conclude 
that $w(a,b) = 1$ when $a$ and $b$ are good. 
\end{proof}

Next, define $\ztilde =  \sum_{a} w(a),$  $\zhat =\sum_{b}w(b),$ and 
$\epsilon_1= 1/6n^2$. We show that $\sum_{a~\text{bad}}w(a)\leq 
\epsilon_1$ and $\sum_{b~\text{bad}}w(b)\leq \epsilon_1/Z_B$.  Thus, we 
find $Z_A\approx \sum_{a~\text{good}}  w(a)$ and $Z_B \approx 
\sum_{b~\text{good}}  w(b)$, and moreover 
$Z=\sum_{a,b}w(a)w(b)w(a,b)\approx Z_AZ_B.$ We show that when $a$ and 
$b$ are both good, $\pi(a)\approx w(a)Z_B/Z$ and $\pi(b)\approx 
w(b)Z_A/Z$.  Thus, when $a$ and $b$ are both good, $r(a,b) = 
\frac{\pi(a,b)}{\pi(a)\pi(b)} \approx \frac{Z}{Z_AZ_B}\approx 1.$

Let $\zAgd = \sum_{a \text{ good}} w(a)$ be the sum $\ztilde$ restricted 
to only good configurations $a$ and $\zBgd = \sum_{b \text{ good}} w(b)$ 
be the sum $\zhat$ restricted to only good configurations $b$. 
By \expref{Lemma}{lem:partitions}, this is a bound on the total weight of 
staircase walks with area more than $N^*$ under the curve.  We will use 
this to bound the weight of the bad $a$'s contributing to $\ztilde$ and 
the weight of the bad $b$'s contributing to $\zhat$.  Specifically, we 
will prove the following lemma.

\begin{lemma}\label{lem:zBound}\ \  
  \begin{enumerate}
  \item $\zAbd =  \sum_{a \text{ bad}} w(a) \leq \epsilon_1$ and $\zBbd =  \sum_{b \text{ bad}} w(b) \leq \epsilon_1\zhat.$
 \item For all $a,b$ we have $\pi(a)\leq w(a)\pi(a^*,b^*)\zhat$ and $\pi(b)\leq w(b)\pi(a^*,b^*)\ztilde.$
 \item For \emph{good} $a$ and \emph{good} $b$ we have $\pi(a) \geq w(a) \pi(a^*,b^*)\zBgd$ and $\pi(b) \geq w(b) \pi(a^*,b^*)\zAgd.$
 \item For all $b$ we have $\pi(b) \geq w(b)\pi(a^*,b^*)w(a^*,b).$
    \end{enumerate}
    \end{lemma}
\begin{proof}
\noindent {\bf Part (1.)} Recall that $w(a) = \pi(a,b^*)/\pi(a^*,b^*).$  Both $k$-particle 
systems $(a,b^*)$ and $(a^*,b^*)$ have the same inversions $(j,l)$ for 
$j,l < i$, so when we consider the ratio $w(a)$, this contains the 
$(i,j)$ for $j > i$ inversions and the $(j,l)$ for $j \geq i, l < i$ 
inversions. Since there are no $(i,j)$ for $j > i$ inversions in 
$(a^*,b^*)$, and the weight of the $(j,l)$ for $j \geq i, l < i$ 
inversions is less in $(a,b^*)$ than it is in $(a^*,b^*)$, we have 
$$\zAbd =  \sum_{a \text{ bad}} w(a) \leq \sum_{a \text{ 
bad}}\prod_{j<l: a(j) > i \atop a(l) = i}q_{a(j), i} \leq \sum_{a 
\text{ bad}}\prod_{j<l: a(j) > i \atop a(l) = i}q_{i+1, i} \leq 
\epsilon_1\,.$$ The last two steps follow from $\p$ being weakly 
monotonic and
from  %
\expref{Lemma}{lem:partitions}, respectively.   

Next, we consider $\zBbd$ and recall that $w(b) = 
\pi(a^*,b)/\pi(a^*,b^*).$  The two $k$-particle systems $(a^*,b)$ and 
$(a^*,b^*)$ have the same inversions $(j,l)$ for $j,l \leq i$.  When 
considering the ratio $w(b)$, there are several types of inversions 
remaining.  There are the $(j,l)$ for $j,l > i+1$ inversions, which 
arise due to the order of the particles in $\mathcal{C}_{>= i+2}$.  We 
will represent these $\tau_{i+2}$.  Additionally, there are the 
inversions $(j,i+1)$ for $j > i+1$, of which there are at least 
$\psize$ if $b$ is \emph{bad}, and the $(j,l)$ for $j \geq i+1, l \leq 
i$ inversions, which are maximized in $\tau_{i+2}^*$, which we will 
define as the highest weight configurations consistent with 
$\sigma_{i-1}$ and $\tau_{i+2}$. In other words, $\tau_{i+2}^*$ is the 
configuration that has the particles in $\mathcal{C}_{< i}$ ordered 
according to $\sigma_{i-1}$, the particles in $\mathcal{C}_{i}$ as far 
forward as possible, then the particles in $\mathcal{C}_{i+1}$ again as 
far forward as possible, and finally the particles in 
$\mathcal{C}_{>i+1}$ in the remaining positions ordered according to 
the $(k-i-1)$-particle system $\tau_{i+2}$.  Given these definitions, 
we have the following:
\begin{align*}
\zBbd &=  \sum_{b \text{ bad}} w(b)=  \sum_{\tau_{i+2}} w (\tau_{i+2}^*) \sum_{EX(i+1,j: j>i+1, \text{ bad})}w(b)/w(\tau_{i+2}^*) \\
&\leq \sum_{\tau_{i+2} }w (\tau_{i+2}^*) \sum_{EX(i+1,j: j>i+1, \text{ bad})}\prod_{j<l: a(j) > i +1\atop a(l) = i+1}q_{a(j), i+1} \\
&\leq \sum_{\tau_{i+2}} w (\tau_{i+2}^*) \sum_{EX(i+1,j: j>i+1, \text{ bad})}\prod_{j<l: a(j) > i +1\atop a(l) = i+1}q_{i+2, i+1} \\
&\leq \sum_{\tau_{i+2}} w (\tau_{i+2}^*)\epsilon_1 \leq \epsilon_1 \zhat,
\end{align*}
where $EX(i+1,j: j>i+1, \text{ bad})$ is an exclusion process 
(2-particle system) with more than $N^*$ inversions; one particle in 
this process is $\mathcal{C}_{i+1}$ and the other is the elements of 
$\mathcal{C}_{>i+1}$ with more than $N^*$ inversions.  In other words, 
we are dividing the $(k-i)$-particle system $b$ based on the location 
of the particles in $\mathcal{C}_{i+1}$.\\


\noindent{\bf Part (2.)}
Both of these bounds are straightforward from the definitions. 
For all $a$ we have the following:
 $$\pi(a) = \sum_{b'} \pi(a,b') = \sum_{b'} w(a)w(b')w(a,b')\pi(a^*,b^*) \leq w(a)\pi(a^*,b^*)\sum_{b'}w(b')= w(a)\pi(a^*,b^*)\zhat\,.$$
Similarly, for all $b$ we have the following:
 $$\pi(b) = \sum_{a'} \pi(b,a') = \sum_{a'} w(a')w(b)w(a',b)\pi(a^*,b^*) \leq w(b)\pi(a^*,b^*)\sum_{a'}w(a')= w(b)\pi(a^*,b^*)\ztilde\,.$$


\noindent{\bf Part (3.)} Next, we give a lower bound on $\pi(a)$ and 
$\pi(b)$ for good $a$ and good $b$.  Using the property that if $a$ and 
$b$ are both good then $w(a,b) = 1$, we have the following bound for 
good $a$:
\begin{align*}
\pi(a) &=  \sum_{b'}\pi(a,b') = \sum_{b'} w(a)w(b') w(a,b')\\
&\geq w(a) \pi(a^*,b^*) \sum_{b' \text{ good}} w(b')w(a,b')\\
&= w(a) \pi(a^*,b^*) \sum_{b' \text{ good}} w(b') = w(a) \pi(a^*,b^*)\zBgd.\\
\end{align*}
\noindent For good $b$ we have a similar lower bound: 
\begin{align*}
\pi(b) &=  \sum_{a'}\pi(a',b) = \sum_{a'} w(a')w(b) w(a',b)\pi(a^*,b^*)\\
&\geq w(b) \pi(a^*,b^*) \sum_{a' \text{ good}} w(a')w(a',b)\\
&= w(b) \pi(a^*,b^*) \sum_{a' \text{ good}} w(a') = w(b) \pi(a^*,b^*)\zAgd.\\
\end{align*}



\noindent{\bf Part (4.)} Finally, we prove a weaker lower bound on $b$ that 
holds for all $b$:
\begin{align*}
\pi(b) &= \sum_{a'}\pi(a',b) = \sum_{a'} w(a')w(b) w(a',b)\pi(a^*,b^*)\\
&= w(b)\pi(a^*,b^*) \sum_a w(a)w(a,b)\\
&\geq w(b)\pi(a^*,b^*) w(a^*)w(a^*,b) = w(b)\pi(a^*,b^*)w(a^*,b).\qedhere
\end{align*}
\end{proof}

Returning to the proof of \expref{Lemma}{lem:perms}, we start by assuming 
both $a$ and $b$ are good.  We will handle the case where either $a$ or $b$ is bad next.  
Recall from \expref{Lemma}{lem:good} that if $a$ and $b$ are both good then $w(a,b) = 1$.  Then, by parts (3) and 
(1) of \expref{Lemma}{lem:zBound}, we have
\begin{align*}
  r(a,b)&=\frac{\pi(a,b)}{\pi(a)\pi(b)}=\frac{w(a)w(b)\pi(a^*,b^*)}{\pi(a)\pi(b)}\\
  &\leq \frac{w(a)w(b)\pi(a^*,b^*)}{(w(a)\pi(a^*,b^*)\zBgd)(w(b) \pi(a^*b^*)\zAgd)}\\
  &= \frac{1}{\pi(a^*,b^*)\zBgd\zAgd}\leq \left(\frac{1}{(1-\epsilon_1)^2}\right) \frac{1}{\zhat\ztilde \pi(a^*,b^*)} \leq \frac{1}{(1-\epsilon_1)^2}.
\end{align*}
The last step uses the following lower bound on $\zhat\ztilde\pi(a^*,b^*)$.
 $$\zhat\ztilde\pi(a^*,b^*) =\sum_{(a,b)} w(a)w(b) \pi(a^*,b^*) > \sum_{(a,b)} w(a)w(b) w(a,b)\pi(a^*,b^*)  = \sum_{(a,b)} \pi(a,b) = 1\,.$$
From part (2) of \expref{Lemma}{lem:zBound} we have the following lower 
bound when $a$ and $b$ are both good:
\begin{align*}
  r(a,b)&=\frac{\pi(a,b)}{\pi(a)\pi(b)}=\frac{w(a)w(b)\pi(a^*,b^*)}{\pi(a)\pi(b)}\\
  &\geq \frac{w(a)w(b)\pi(a^*,b^*)}{
     \left(w(a)\pi(a^*,b^*)\zhat\right)\left(w(b)\pi(a^*,b^*)\ztilde\right)}\\
  &= \frac{1}{
     \zhat\ztilde\pi(a^*,b^*)} \geq(1-\epsilon_1)^2.
\end{align*}
In the last step we upper bound $\zhat\ztilde\pi(a^*,b^*)$ as follows 
using part (1) of \expref{Lemma}{lem:zBound}:
\begin{equation*}\label{prodBound}
\zhat\ztilde\pi(a^*,b^*) \leq\frac{\zBgd\zAgd\pi(a^*,b^*)}{ (1-\epsilon_1)^{2}}= \sum_{a~\text{good},\atop b~\text{good}} \frac{w(a)w(b) \pi(a^*,b^*)}{(1-\epsilon_1)^{2}}< \sum_{(a,b)}\frac{ \pi(a,b)}{(1-\epsilon_1)^{2}} = \frac{1}{(1-\epsilon_1)^{2}}.
\end{equation*}
In order to apply \expref{Lemma}{lem:directproduct}, we need to bound the 
quantity $\sum_{(a,b)}\pi(a,b)\left(\sqrt{r(a,b)}-\frac{1}{\sqrt{r(a,b)}}\right)^2$.  
There are two cases depending on whether $r(a,b)\leq 1$, but either
way we have
\[\left(\sqrt{r(a,b)}-\frac{1}{\sqrt{r(a,b)}}\right)^2\leq
  \left(\frac{1}{1-\epsilon_1} - (1-\epsilon_1)\right)^2\leq
  5\epsilon_1^2\,,\]
as long as $\epsilon_1\leq .191$ (this is true since $\epsilon_1\leq 
1/(6n^2)$ and $n\geq 2$).  Therefore
  \[\sum_{a~\text{good},\atop
    b~\text{good}}\pi(a,b)\left(\sqrt{r(a,b)}-\frac{1}{\sqrt{r(a,b)}}\right)^2\leq
  5\epsilon_1^2\,.\]
  
Next, we address the case where at least one of $a$ or $b$ is bad.  In this case, we
will show that the weight of these configurations is so small that
it overcomes the fact that $w(a,b)$ and $r(a,b)$ may be exponentially small.  
If $r(a,b)\leq 1$ then
$\pi(a,b)\left(\sqrt{r(a,b)}-\frac{1}{\sqrt{r(a,b)}}\right)^2\leq
\pi(a)\pi(b)$.  Otherwise,
$\pi(a,b)\left(\sqrt{r(a,b)}-\frac{1}{\sqrt{r(a,b)}}\right)^2\leq
\frac{\pi(a,b)^2}{\pi(a)\pi(b)}.$  Either way,
\[\pi(a,b)\left(\sqrt{r(a,b)}-\frac{1}{\sqrt{r(a,b)}}\right)^2\leq
\pi(a)\pi(b) + \frac{\pi(a,b)^2}{\pi(a)\pi(b)}\,.\]
In order to upper bound $\frac{\pi(a,b)^2}{\pi(a)\pi(b)}$, we will
bound the conditional probabilities $\pi(a,b)/\pi(a)$ for \emph{good} 
$a$ and $\pi(a,b)/\pi(b)$ for \emph{all} $b$.  Using part (3) of \expref{Lemma}{lem:zBound} 
we have for $a$ good:
$$\frac{\pi(a,b)}{\pi(a)} \leq \frac{w(a)w(b)w(a,b)\pi(a^*,b^*)}{w(a)\pi(a^*,b^*)\zBgd} \leq \frac{w(b)}{\zBgd}\,.$$
Using part (4) of \expref{Lemma}{lem:zBound} we have for all $b$:
$$\frac{\pi(a,b)}{\pi(b)} \leq \frac{w(a)w(b)w(a,b)\pi(a^*,b^*)}{w(b)w(a^*,b)\pi(a^*,b^*)} \leq w(a)\,.$$
 
\noindent Now we may use these facts to bound $\sum_{a,b} 
\frac{\pi(a,b)^2}{\pi(a)\pi(b)}$ for \emph{bad} $a$ and \emph{all} $b$ 
using part (1) of \expref{Lemma}{lem:zBound}:
 \begin{align*}
   \sum_{a \text{ bad},b} \frac{\pi(a,b)^2}{\pi(a)\pi(b)}
      &\leq  \sum_{a \text{ bad},b} \frac{\pi(a,b)}{\pi(a)}w(a) \\ 
   &= \sum_{a \text{ bad}} w(a) \frac{ \sum_b \pi(a,b)}{\pi(a)} \\ 
 &= \sum_{a \text{ bad}} w(a)~~ \leq ~~\epsilon_1.
 \end{align*}
 
\noindent Similarly, we bound $\sum_{a,b} \frac{\pi(a,b)^2}{\pi(a)\pi(b)}$ 
for good $a$ and bad $b$ using part (1) of \expref{Lemma}{lem:zBound}.
  \begin{align*}
    \sum_{a \text{ good},b\text{ bad}} \frac{\pi(a,b)^2}{\pi(a)\pi(b)}
    &\leq \sum_{a \text{ good},b\text{ bad}}
    \frac{\pi(a,b)}{\pi(b)}\frac{w(b)}{\zBgd}\\
    &\leq \sum_{b\text{ bad}} \frac{w(b)}{\zBgd}\\
    &= \frac{\zBbd}{\zBgd}\\
    &\leq \frac{\epsilon_1}{1- \epsilon_1}.
 \end{align*}
 
We continue our analysis of the case where at least one of $a$ or $b$ is bad and bound $\sum_{a,b}\pi(a)\pi(b).$ 
Notice that the summations over $a$ and $b$ are separable, so we have 
the following:
\begin{equation*}
  \sum_{a~\text{bad}, b} \pi(a)\pi(b) = \Pr(a~\text{bad})
\end{equation*}
and
\begin{equation*}
  \sum_{b~\text{bad}, a} \pi(a)\pi(b) = \Pr(b~\text{bad}).
\end{equation*}
We may bound these probabilities via the following inequalities.
 $$\zhat\pi(a^*,b^*) \leq
  \frac{\zBgd\zAgd\pi(a^*,b^*)}{(1-\epsilon_1)}=\frac{1}{(1-\epsilon_1)}\sum_{a~\text{good},\atop
    b~\text{good}} w(a)w(b) \pi(a^*,b^*)  <
  \frac{1}{1-\epsilon_1}\,.$$
  
\noindent Similarly, we have $\ztilde\pi(a^*,b^*)\leq 
1/(1-\epsilon_1)$.  Using part (2) of \expref{Lemma}{lem:zBound}, this 
implies that
  \[\Pr(a~\text{bad}) = \sum_{a~\text{bad}}\pi(a) \leq
  \sum_{a~\text{bad}}w(a)\pi(a^*,b^*)\zhat\leq
  \sum_{a~\text{bad}}\frac{w(a)}{1-\epsilon_1}\leq
  \frac{\epsilon_1}{1-\epsilon_1} \,.\]
  We also have $\Pr(b~\text{bad}) \leq \epsilon_1/(1-\epsilon_1)$.
Thus when at least one of $a$ and $b$ is bad we have $$\sum_{a,b}\pi(a,b)\left(\sqrt{r(a,b)}-\frac{1}{\sqrt{r(a,b)}}\right)^2\leq
\sum_{a,b} \left(\pi(a)\pi(b) + \frac{\pi(a,b)^2}{\pi(a)\pi(b)}\right) \leq \epsilon_1 + \frac{3\epsilon_1}{1-\epsilon_1}\,.$$

Putting both cases together, we have
 \[\sum_{a, b}\pi(a,b)\left(\sqrt{r(a,b)}-\frac{1}{\sqrt{r(a,b)}}\right)^2\leq
  5\epsilon_1^2 + \epsilon_1 + \frac{3\epsilon_1}{1-\epsilon_1} \leq
  6\epsilon_1 = 1/n^2\,,\]
  as long as $\epsilon_1\leq .225$, which is true since $n\geq 2$.
Thus, this decomposition is %
$\epsilon$-orthogonal, where $\epsilon\leq 1/n^2$ %
This completes the proof of \expref{Lemma}{lem:perms}. %
%
%
%
\end{proof}   %

\subsection{The Markov chain $\mn$}\label{sec:mn}

Here we give the details to show how we can use \expref{Theorem}{pprocess}, 
which gives a  %
lower  %
bound on the spectral gap of the particle process Markov 
chain $\mk$, to obtain %
an upper  %
bound on the mixing time of the 
nearest-neighbor Markov chain $\mn.$  

First, we begin with some preliminaries on Markov chains and mixing 
times.  The time a Markov chain takes to converge to its stationary 
distribution, or \emph{mixing time}, is measured in terms of the 
distance between the distribution at time~$t$ and the stationary 
distribution $\pi$. %
 The \emph{total variation distance} at time~$t$ is 
$\|P^t,\pi\| _{tv} = \max_{x\in\Omega}\frac{1}{2}\sum_{y\in\Omega} 
|P^t(x,y)-\pi(y)|,$  %
where $P^t(x,y)$ is the $t$-step transition 
probability, 
and $\Omega$ is the state space of the Markov chain.  %
For all $\epsilon>0$, the \emph{mixing time} 
$\tau(\epsilon)$ of $\m$ is defined as
$$\tau(\epsilon)=\min \{t: \|P^{t'},\pi \|_{tv}\leq \epsilon, \forall t' \geq t\}\,.$$ 

In order to use the lower bound on the spectral gap to obtain an upper bound on the mixing time we will use the following well-known result. 
Again, let $\Gap(\transmat)=1-|\lambda_1|$ denote the spectral gap, where $\lambda_0, \lambda_1, \ldots, \lambda_{|\Omega|-1}$ are the eigenvalues of the transition matrix $\transmat$ and $1=\lambda_0>|\lambda_1|\geq |\lambda_i|$ for all $i\geq 2$. The following result  relates the spectral gap with the mixing time (see, \eg, \cite{sinclair},\cite{RM03}):

%
\begin{theorem}[\cite{RM03}]\label{gap}
Let $\pi_*=\min_{x\in\Omega} \pi (x)$. For all $\epsilon>0$ we have
\begin{align*}
(a)\qquad \tau(\epsilon)&\leq \frac{1}{\Gap(\transmat)} \log\left(\frac{1}{\pi_*\epsilon}\right).\\. %
(b)\qquad \tau(\epsilon)&\geq \frac{|\lambda_1|}{2\Gap(\transmat)} \log\left(\frac{1}{2\epsilon}\right).
%
%
\end{align*}
\end{theorem}
%
%

We also require the following comparison theorem
%
%
{}from~\cite{dsc}.   %
Let $P'$ and $P$ be two
reversible Markov chains on the same state space $\Omega$ with the same 
stationary distribution $\pi$ and let $E(P) = \{(x,y): P(x,y) >
0\}$ and $E(P') = \{(x,y): P'(x,y) > 0\}$ denote the sets of edges of
the two chains, viewed as directed graphs.  For each $x,y$ with
$P'(x,y)>0$, define a path $\gamma_{xy}$ using a sequence of states
$x=x_0,x_1,\cdots,x_k = y$ with $P(x_i,x_{i+1})>0$, and let
$|\gamma_{xy}|$ denote the length of the path.  Let $\Gamma(z,w) =
\{(x,y) \in E(P'): (z,w) \in \gamma_{xy}\}$ be the set of paths that
use the transition $(z,w)$ of $P$.  Finally, define  
\abovedisplayskip=3pt
\belowdisplayskip=3pt
\begin{equation}\label{eqn:zeta}
\zeta = \max_{(z,w) \in E(P)} \left \{\frac{1}{\pi(z)P(z,w)}
\sum_{\Gamma(z,w)}|\gamma_{xy}|\pi(x)P'(x,y) \right \}.
\end{equation}

\begin{theorem}[\cite{dsc}]\label{thm_comp}
With the above notation, $\Gap(P)\geq \frac{1}{\zeta} \Gap(P')$.
\end{theorem}

Now we may prove \expref{Theorem}{thm:mnGap}.\\

%
\begin{proof}[Proof of \expref{Theorem}{thm:mnGap}]
Given our improved bound on the spectral gap of $\mk$ from
\expref{Theorem}{pprocess}, the remainder of this proof uses exactly the
same approach as~\cite{MS}, except that we use \expref{Corollary}{cor:product}
to eliminate one factor of $n$.  We include a summary here for
completeness.  The complete details can be found in~\cite{MS}.
Instead of analyzing $\mn$ directly we will analyze an auxiliary chain
$\mtk$ that allows a larger set of transpositions (including those
allowed by $\mk$) and then use the comparison theorem, \expref{Theorem}{thm_comp},
to obtain a bound for $\mn$.   
In what follows, we write $C(x)$ to mean the index of the class containing element $x\in [n]$, so if $x\in C_i$ then $C(x)=i$.

\vspace{.1in}
\noindent  {\bf The Transposition Markov chain $\mtk$} 

\vspace{.05in}
\noindent {\tt Starting at any permutation $\sigma^0$, iterate the following:} 
\begin{itemize}
\item At time $t\geq0,$ choose a position $1\leq i \leq n$ in permutation $\sigma^t$ and a direction $d \in \{L, R,N\}$ uniformly at random.  
\item If $d = L$, find the largest $j$ with $1 \leq j<i$ and $C(\sigma^t(j)) \geq C(\sigma^t(i))$ (if one exists).  If $C(\sigma^t(j)) > C(\sigma^t(i)),$ then 
with probability $1/2$,
exchange $\sigma^t(i)$ and $\sigma^t(j)$ to obtain $\sigma^{t+1}.$  
\item If $d = R$, find the smallest $j$ with $n\geq j>i$ and $C(\sigma^t(j)) \geq C(\sigma^t(i))$ (if one exists).  If $C(\sigma^t(j)) > C(\sigma^t(i)),$ then with probability 
$$\frac{1}{2}~ q_{\sigma^t(j),\sigma^t(i)}  \prod_{i<k<j}\left(q_{\sigma^t(j),\sigma^t(k)}q_{\sigma^t(k),\sigma^t(i)}\right)\,,$$
exchange $\sigma^t(i)$ and $\sigma^t(j)$ to obtain $\sigma^{t+1}.$  
\item If $d = N,$  find the largest $j$ with $1\leq j<i$ and $C(\sigma^t(j)) = C(\sigma^t(i)).$  If such an element exists, then with probability $1/2$, exchange the elements $\sigma^t(i)$ and $\sigma^t(j)$ to obtain $\sigma^{t+1}.$ 

\item With all remaining probability, $\sigma^{t+1} = \sigma^t.$
\end{itemize}

The Markov chain $\mtk$ has the same
stationary distribution as $\mn$ and is a product of
$k+1$ independent Markov chains~\cite{MS}.  The first~$k$ chains involve moves
between particles in the same class and the $i$th is an unbiased
nearest-neighbor Markov chain over permutations on $c_i$ particles.
%
%
%
This chain has spectral gap $\Theta(c_i^{-3})$. (This is an 
unpublished result of Diaconis. See, \eg,~\cite{wilson}.)    %
\begin{lemma}[Diaconis]
The spectral gap of the unbiased nearest neighbor Markov chain
over permutations on $[n]$ is $(1-\cos(\pi/n))/(n-1)=\Theta(n^{-3})$.
\end{lemma}
\noindent However, the chain $\mtk$ updates one of the first $k$ chains
only if direction $N$ is selected,  which happens with probability $c_i/(6n)$
for each particle class $1\leq i \leq k.$  Thus, the spectral gap of
the slowed-down version of this chain is $\Theta(1/(nc_i^2))$.
The final chain $\mk$ which we analyzed in \expref{Section}{sec:perm}
allows only moves between different particle classes; it is updated when
direction $L$ or $R$ is selected (\ie,  with probability $2/3$),
so,  %
by
\expref{Theorem}{pprocess}, it has spectral gap  $\Omega(n^{-2})$.
Therefore, by \expref{Corollary}{cor:product}, the spectral gap of $\mtk$
is $\Omega(n^{-3})$.

The final step is to relate the spectral gap of $\mn$ to that of $\mtk$
using the comparison theorem of~\cite{dsc}, \expref{Theorem}{thm_comp}.  
Section~5 of~\cite{MS} proves the following lemma where $\zeta$ is 
defined as in \expeqref{equation}{eqn:zeta}.

\begin{lemma}[\cite{MS}]\label{lem_comp} If the probabilities~$\p$ are
  weakly monotonic and form a bounded \kP\, for $k\geq 2$, then $\zeta=O(n^4)$.
\end{lemma}

\noindent Combining \expref{Lemma}{lem_comp} and \expref{Theorem}{thm_comp}
with our bound on the gap of $\mtk$, we get $\Gap(\mn)=\Omega(n^{-7})$. 
%
\end{proof}

\vspace{.2in}

Finally, we may bound the mixing time of $\mn$. 

\begin{theorem}\label{mnMixing} If the %
%
probability array $\cP$ is
weakly monotonic and
%
forms  %
a bounded \kP\, for $k\geq 2$, with $|C_i|\geq 2\psize$, 
then the mixing time $\tau_{\text{nn}}$ of $\mn$ satisfies 
$\tau_{\text{nn}} = O(n^{9}\ln(1/\epsilon))$.
\end{theorem}
\begin{proof}
Finally, to get the mixing time of $\mn$, we let $q_* =
\max_{i<j} q_{i,j}$ then $\pi_* =  \min_{x\in \Omega}\pin(x)
\geq (q_*^{\binom{n}{2}}n!)^{-1}$ (see \cite{bmrs} and \cite{MS} for
more details), so $\log(1/\epsilon \pi_*) = O(n^2\ln \epsilon^{-1})$
since $q$ is bounded from above by a positive constant (all $p_{i,j}$ are constant with respect to $n$).  Applying
\expref{Theorem}{gap}(a) we have that the
mixing time of $\mn$ is $O(n^9\ln\epsilon^{-1})$.  This proves
\expref{Theorem}{mnMixing}.
\end{proof}
